\chapter{Empirical Evaluation}
\label{ch:hitl}

This chapter answers the three research questions. It opens with the system-level result: the full-pipeline ablation (Section~\ref{sec:pipeline_ablation}) that isolates what every stage contributes on every canonical surface, under the single protocol of Section~\ref{sec:eval_protocol}. Temporal (segment-level) evidence follows on the consecutive-frame video surface (Section~\ref{sec:temporal_results}). The component deep-dives (Section~\ref{sec:components_eval}) then examine each stage (the RGB detector's three training stances, the IR detector's six-revision HITL case study, the trust-classifier comparison, and the filter comparison), including the design-evolution evidence measured under earlier configurations, always labelled as such. The Model MRI findings (Section~\ref{sec:mri_findings}) and the grayscale cross-modal finding (Section~\ref{sec:grayscale}) close the results; threats to validity and limitations close the chapter. Every table states its $n$; headline cells carry $95\%$ bootstrap confidence intervals in brackets. Such an interval is the central $95\%$ range of the metric over $1{,}000$ resamples of the evaluation frames drawn with replacement: it indicates how much the number would move under a different sample of frames, so two cells whose intervals do not overlap differ by more than sampling noise.

% ======================================================================
\section{Full-Pipeline Ablation}
\label{sec:pipeline_ablation}
\label{ch:experiments}

\paragraph{Stack under test.}
The production stack: RGB detector \texttt{ft4}, IR detector \texttt{v3b}, trust router \texttt{robust8} ($\tau=0.20$), and the two per-frame confuser filters. Two different threshold stages should not be conflated. A \emph{detector confidence floor} discards a detection before it enters the pipeline; these caches run both detectors at \texttt{conf}$=$0.25 (the routing classifier's training setting). The deployment IR floor is $0.40$, but because the filter rescores every surviving detection the evaluation is insensitive to that floor and uses one value for both detectors (Section~\ref{sec:design_rationale}). A \emph{filter threshold} is then applied to the filter's own $P(\text{drone})$ output, the probability above which a surviving detection is accepted as a drone: \texttt{mlp\_v5} (RGB) at $0.25$, and \texttt{mlp\_v5\_ir\_aligned} at $0.05$ under its thermal input scaler and $0.25$ under its grayscale scaler. The IR filter's $0.05$ is therefore a filter threshold on thermal input, not the IR detector's deployment floor. Ablation arms swap exactly one element at a time: the router (\texttt{sa32}, the strongest hand-engineered classifier; \texttt{robust6}, the six-feature statistical base), the filter (the superseded MobileNetV3 patch CNN at \texttt{patch\_thr}$=$0.5), the composition order (classifier$\to$filter vs filter$\to$classifier), the router's reject discipline (a single global reject-probability floor, the \texttt{rej$\geq$0.8} rows), or the stage's presence. All cells replay the same cached detections (Section~\ref{sec:unified_cache}); scoring is per-modality trust-aware (Section~\ref{sec:scoring_audit}).

\subsection{Paired Drone Benchmarks: the Cost Side}
\label{sec:pipeline_paired}

On drone-positive paired surfaces the question is what the pipeline \emph{costs}: every vetoed detection here is potentially a missed drone.

\begin{table}[htbp]
    \centering
    \caption{Full-pipeline ablation, Svanstr\"om paired ($n=4{,}000$ frames, even-strided across the full 28{,}710-frame set, not the first $4{,}000$; RGB @1280 / IR @640, IoP@0.5). \emph{RGB only}/\emph{IR only} score that detector alone against its own ground truth; every row below the second rule is trust-aware and \emph{sums both modalities} (RGB-side scored vs RGB GT $+$ IR-side vs IR GT), so its TP/FP/FN base is both modalities' drones combined, not the frame count. \texttt{clf} $=$ $+$trust router; \texttt{filt} $=$ $+$per-frame filter; \texttt{clf$\to$filt} $=$ router then filter. The shipped production cell is \texttt{filt$\to$clf[robust8-nr]} (filter then the no-reject router; \texttt{robust8} with the reject class is the adjacent ablation, Section~\ref{sec:classifier_results}); the pipeline \emph{raises} recall over bare ($0.948\to0.991$) while adding $\sim$29~pp precision. \texttt{clf$\to$filt} is the same router in the other order (Section~\ref{sec:classifier_results} on why \texttt{filt$\to$clf} is shipped).}
    \label{tab:ablation_svanstrom}
    \begin{tabular}{lcccccl}
        \toprule
        cell & TP & FP & FN & $P$ & $R$ & F1 [95\% CI] \\
        \midrule
        RGB only (\texttt{ft4}, vs RGB GT)   & 1532 & 1845 & 141 & 0.454 & 0.916 & 0.607 [0.590--0.623] \\
        IR only (\texttt{v3b}, vs IR GT)     & 1610 &  174 &  31 & 0.903 & 0.981 & 0.940 [0.931--0.949] \\
        \midrule
        bare (both detectors)        & 3142 & 2019 & 172 & 0.609 & 0.948 & 0.742 [0.730--0.754] \\
        filt only (\texttt{mlp\_v5}, RGB)    & 3043 &  353 & 271 & 0.896 & 0.918 & 0.907 [0.900--0.914] \\
        filt only (\texttt{aligned}, IR)     & 3142 & 2018 & 172 & 0.609 & 0.948 & 0.742 [0.730--0.754] \\
        filt only (patch)            & 3004 &  782 & 310 & 0.793 & 0.907 & 0.846 [0.836--0.855] \\
        clf only [robust8-nr]        & 3136 & 1997 &  27 & 0.611 & 0.992 & 0.756 [0.745--0.768] \\
        clf only [robust8]           & 3092 &  354 &  31 & 0.897 & 0.990 & 0.941 [0.935--0.947] \\
        clf only [robust6]           & 3132 &  286 &  34 & 0.916 & 0.989 & 0.951 [0.946--0.957] \\
        clf only [sa32]              & 3137 &  187 &  30 & 0.944 & 0.991 & 0.967 [0.962--0.971] \\
        clf$\to$filt [robust8-nr]    & 3037 & 330 & 126 & 0.902 & 0.960 & 0.930 [0.923--0.937] \\
        clf$\to$filt [robust8]       & 2993 & 195 & 130 & 0.939 & 0.958 & 0.949 [0.943--0.954] \\
        clf$\to$filt [sa32]          & 3039 &  162 & 128 & 0.949 & 0.960 & 0.955 [0.949--0.960] \\
        \textbf{filt$\to$clf [robust8-nr]} & 3037 & 332 & 29 & \textbf{0.901} & \textbf{0.991} & \textbf{0.944 [0.938--0.950]} \\
        filt$\to$clf [robust8]       & 2996 &  196 &  35 & 0.939 & 0.989 & 0.963 [0.957--0.968] \\
        \bottomrule
    \end{tabular}
\end{table}
% [source: thesis_eval/results/tier1_results.json (svanstrom B_pipeline); run=thesis_eval/pipeline_eval_unified.py]

\begin{table}[htbp]
    \centering
    \caption{Full-pipeline ablation, Anti-UAV paired ($n=4{,}000$ frames, even-strided across the 85{,}374-frame set, both modalities @640, IoU@0.5). Row conventions as in Table~\ref{tab:ablation_svanstrom} (\emph{RGB/IR only} score one detector vs its own GT; the trust-aware rows sum both modalities). The no-harm control: every pipeline cell is within $\sim$1~pp of bare or above it, the shipped no-reject \texttt{robust8-nr} included (this saturated surface rarely rejects, so dropping the reject class changes nothing here).}
    \label{tab:ablation_antiuav}
    \begin{tabular}{lcccccl}
        \toprule
        cell & TP & FP & FN & $P$ & $R$ & F1 [95\% CI] \\
        \midrule
        RGB only (\texttt{ft4}, vs RGB GT)   & 3657 &  41 &  68 & 0.989 & 0.982 & 0.985 [0.982--0.988] \\
        IR only (\texttt{v3b}, vs IR GT)     & 3775 & 133 & 173 & 0.966 & 0.956 & 0.961 [0.956--0.966] \\
        \midrule
        bare (both detectors)        & 7432 & 174 & 241 & 0.977 & 0.969 & 0.973 [0.970--0.976] \\
        filt only (\texttt{mlp\_v5}, RGB)    & 7430 & 172 & 243 & 0.977 & 0.968 & 0.973 [0.970--0.976] \\
        filt only (\texttt{aligned}, IR)     & 7432 & 173 & 241 & 0.977 & 0.969 & 0.973 [0.970--0.976] \\
        filt only (patch)            & 7430 & 174 & 243 & 0.977 & 0.968 & 0.973 [0.970--0.976] \\
        clf only [robust8-nr]        & 7430 & 160 &  81 & 0.979 & 0.989 & 0.984 [0.981--0.986] \\
        clf only [robust8]           & 7408 & 141 &  93 & 0.981 & 0.988 & 0.984 [0.982--0.987] \\
        clf only [robust6]           & 7384 & 156 & 170 & 0.979 & 0.978 & 0.978 [0.975--0.981] \\
        clf only [sa32]              & 7426 & 135 &  75 & 0.982 & 0.990 & 0.986 [0.984--0.988] \\
        clf$\to$filt [robust8-nr]    & 7428 & 157 & 83 & 0.979 & 0.989 & 0.984 [0.982--0.987] \\
        clf$\to$filt [robust8]       & 7406 & 139 & 95 & 0.982 & 0.987 & 0.984 [0.982--0.987] \\
        clf$\to$filt [sa32]          & 7424 & 134 &  77 & 0.982 & 0.990 & 0.986 [0.984--0.988] \\
        \textbf{filt$\to$clf [robust8-nr]} & 7428 & 157 & 81 & \textbf{0.979} & \textbf{0.989} & \textbf{0.984 [0.982--0.987]} \\
        filt$\to$clf [robust8]       & 7406 & 139 &  93 & 0.982 & 0.988 & 0.985 [0.982--0.987] \\
        \bottomrule
    \end{tabular}
\end{table}
% [source: thesis_eval/results/tier1_results.json (antiuav B_pipeline); run=thesis_eval/pipeline_eval_unified.py]

Five readings, each attributable to one element.

\emph{The pipeline maintains recall, and in fact raises it (RQ1).} On Svanstr\"om the production cell's recall is \emph{higher} than bare ($0.948 \to 0.991$) while precision rises $0.609 \to 0.901$. That precision gain is RGB-side: the bare detector's confuser false positives (the RGB-only row alone carries $1{,}845$ of them) are removed by the \texttt{mlp\_v5} filter, whose drone/confuser separability is established independently by the Model MRI (Section~\ref{sec:mri_findings}). The recall gain is a separate, routing effect under per-modality scoring: the router stops the system being scored on the modality that is blind on a given frame, so its misses leave the books (bare FN $172 \to 27$ at the router stage). No new detection is created downstream; the architecture cannot do that. On Anti-UAV the pipeline is no-harm by every cell, which is exactly the role of a saturated control.

\emph{In the shipped stack the filter suppresses the false positives; the reject class is a complementary net we trade away (RQ2).} The shipped no-reject router routes every frame, so on Svanstr\"om it removes almost no bare false positives on its own ($2019 \to 1997$); the \texttt{mlp\_v5} filter does that work ($2019 \to 353$ alone, $\to 330$ composed). The reject-class \texttt{robust8} ablation tells the complementary half of the story: its router \emph{alone} removes $82\%$ of bare FPs ($2019 \to 354$) by vetoing whole untrustworthy frames, a different failure mode than the filter's per-detection veto, which is why composing \texttt{robust8} with the filter reaches $195$ FP. Each stage catches what the other misses: the filter cannot recover a frame whose trusted modality is wrong, and the reject class cannot veto a single confuser detection inside a trusted frame. We ship the no-reject router and delegate false-positive removal to the filter (Section~\ref{sec:classifier_results}); \texttt{robust8} keeps that second net at the recall cost quantified throughout this section.

\emph{Composition order is a recall/precision dial, and for the no-reject router \texttt{filt$\to$clf} is the better side of it.} \texttt{filt$\to$clf} keeps recall at the router's level ($0.991$ on Svanstr\"om) because the router re-routes on the already-filtered detections; \texttt{clf$\to$filt} trades $\sim$3~pp of that recall ($0.960$) for marginal precision. With a reject class the order mattered for confuser suppression (reject vetoes whole frames \emph{before} the filter, the historic reason \texttt{clf$\to$filt} was shipped), but the no-reject router never vetoes a confuser frame, so the two orders give \emph{identical} confuser fire (Table~\ref{tab:ablation_confusers}); with that tie broken we ship \texttt{filt$\to$clf} for the recall (Svanstr\"om F1 $0.944$ vs $0.930$, DUT $0.837$ vs $0.792$).

\emph{Why the no-reject router is shipped.} Among the reject-class routers, the 32-feature \texttt{sa32} leads every in-domain cell by $0.2$--$1.8$~pp, but it needs an expensive per-frame OpenCV pass, was trained against a superseded detector's calibration, and leaks scene statistics (Section~\ref{sec:feature_selection}); \texttt{robust8} matches it within a point on eight free features at a $404\times$ cheaper decision, with \texttt{robust6} between them. On the confuser-rich Svanstr\"om flagship the reject class earns its keep (at the shipped \texttt{filt$\to$clf} order, $\texttt{robust8}$ $0.963$ vs the no-reject $0.944$, at far lower confuser fire). The shipped router is nonetheless the no-reject \texttt{robust8-nr}: a single global router must serve paired, single-modality, and degraded-channel deployments alike, and the reject class collapses recall on the latter two (Tables~\ref{tab:ablation_dut},~\ref{tab:ablation_solo}); averaged across surfaces the no-reject router leads ($0.834$ vs $0.733$ mean composed F1). \texttt{robust8}/\texttt{robust6}/\texttt{sa32} are retained as ablations and the trade is detailed in Section~\ref{sec:classifier_results}.

\emph{Dual-modality routing earns its keep on coverage (RQ3).} Table~\ref{tab:rq3} scores each modality against its own GT and the routed system trust-aware on the same frames. The per-surface reversal is the point: IR dominates Svanstr\"om ($0.940$ vs RGB's $0.607$), RGB dominates Anti-UAV ($0.985$ vs IR's $0.961$); a fixed single-modality choice loses badly somewhere. Under trust-aware scoring the routed no-reject system tracks or edges the stronger modality on each surface (Svanstr\"om $0.944$, just above IR-only's $0.940$; Anti-UAV $0.984$), because that rule deliberately does not charge a modality for what it cannot see (Section~\ref{sec:scoring_audit}). The decisive RQ3 evidence is therefore the \emph{coverage} comparison (Table~\ref{tab:modality_ab}), where every modality's ground truth always counts: there the routed system reaches $0.909$ on Svanstr\"om against RGB-only $0.582$ and IR-only $0.632$, the coverage a single camera physically cannot match.

\begin{table}[htbp]
    \centering
    \caption{RQ3: single modality vs routed pipeline (bare per-modality rows vs the production cell). Each modality is scored against its own GT; the routed row is trust-aware. The modality ranking \emph{reverses} between surfaces; under trust-aware scoring the routed no-reject system tracks the stronger modality (the fair, coverage-scored comparison that isolates routing's value is Table~\ref{tab:modality_ab}).}
    \label{tab:rq3}
    \begin{tabular}{lll}
        \toprule
        configuration & Svanstr\"om F1 & Anti-UAV F1 \\
        \midrule
        RGB only (\texttt{ft4}, bare)  & 0.607 [0.590--0.623] & 0.985 [0.982--0.988] \\
        IR only (\texttt{v3b}, bare)   & 0.940 [0.931--0.949] & 0.961 [0.956--0.966] \\
        \textbf{Routed (production, \texttt{robust8-nr}, filt$\to$clf)}   & \textbf{0.944 [0.938--0.950]} & \textbf{0.984 [0.982--0.987]} \\
        \bottomrule
    \end{tabular}
\end{table}
% [source: thesis_eval/results/tier1_results.json (A_bare + B_pipeline)]

\paragraph{A clean in-domain check: DUT Anti-UAV (\emph{test split}).}
The two surfaces above are IR-dominant (Svanstr\"om) or saturated (Anti-UAV). The DUT Anti-UAV \emph{test split}~\cite{zhao2022dutantiuav} adds an RGB-strong, confuser-free check: an all-drone RGB benchmark on which the RGB detector \texttt{ft4} is in-domain (DUT was one source in the composite \texttt{rgb\_dataset} it trained on), paired with the thermal \texttt{v3b} detector fed grayscale-converted RGB, i.e.\ the no-thermal cross-domain fallback. \textbf{The grayscale-IR column measures this fallback mode, not the IR detector's native thermal performance}: on real thermal it reaches $F1\,0.940$ (Svanstr\"om, Table~\ref{tab:rq3}) and $0.961$ (Anti-UAV); the much lower numbers here reflect a thermal model run on grayscale RGB, an input it was never trained on, and should not be read as the IR detector's capability. Only the official test split is used (no DUT frame here was a training image for the router or filters). Because no confuser is present, every filter veto can only remove a true positive, so the pipeline's one available contribution is precision. We report at \texttt{imgsz}$=$960, the resolution-optimal setting for this small-drone surface (Section~\ref{sec:resolution_arch}; drones down to $25\times9$~px need the pixels).

\begin{table}[htbp]
    \centering
    \caption{Full-pipeline ablation, DUT Anti-UAV \emph{test split} ($n=2{,}200$ frames, $2{,}245$ GT drones, RGB \texttt{ft4} $+$ grayscale-IR \texttt{v3b}, both @960, IoU@0.5, trust-aware). The grayscale-IR channel is the IR detector's no-thermal \emph{fallback} mode (thermal model on grayscale RGB), not its native thermal performance. The full test split is used (not strided). Row conventions as in Table~\ref{tab:ablation_svanstrom} (\emph{RGB/IR only} score one detector vs its own GT; the trust-aware rows sum both modalities). Confuser-free (all-drone): filters can only delete true positives, so here the pipeline buys precision, not FP suppression. \texttt{filt$\to$clf[robust8-nr]} (the shipped no-reject router) is the production configuration; dropping the reject class recovers the recall \texttt{robust8} surrenders on this clean surface (at \texttt{filt$\to$clf}, $R\,0.583\to0.808$, F1 $0.707\to0.837$).}
    \label{tab:ablation_dut}
    \begin{tabular}{lcccccl}
        \toprule
        cell & TP & FP & FN & $P$ & $R$ & F1 [95\% CI] \\
        \midrule
        RGB only (\texttt{ft4}, vs RGB GT)            & 1917 & 104 &  328 & 0.949 & 0.854 & 0.899 [0.889--0.908] \\
        grayscale-IR only (\texttt{v3b}, vs IR GT)    & 1113 & 375 & 1132 & 0.748 & 0.496 & 0.596 [0.576--0.616] \\
        \midrule
        RGB \texttt{ft4} $+$ grayscale-IR \texttt{v3b} & 3030 & 479 & 1460 & 0.864 & 0.675 & 0.758 [0.747--0.769] \\
        filt only (\texttt{mlp\_v5}, RGB)        & 2820 & 455 & 1670 & 0.861 & 0.628 & 0.726 [0.715--0.739] \\
        filt only (\texttt{aligned\_gray}, IR)   & 2796 & 423 & 1694 & 0.869 & 0.623 & 0.725 [0.715--0.736] \\
        filt only (patch)            & 2931 & 470 & 1559 & 0.862 & 0.653 & 0.743 [0.732--0.754] \\
        clf only [robust8-nr]        & 3024 & 464 &  522 & 0.867 & 0.853 & 0.860 [0.850--0.869] \\
        clf only [robust8]           & 2586 & 302 & 1301 & 0.895 & 0.665 & 0.763 [0.749--0.778] \\
        clf only [robust6]           & 2415 & 374 & 1808 & 0.866 & 0.572 & 0.689 [0.674--0.705] \\
        clf only [sa32]              & 2651 & 430 & 1480 & 0.860 & 0.642 & 0.735 [0.721--0.750] \\
        clf$\to$filt [robust8-nr]    & 2580 & 389 & 966 & 0.869 & 0.728 & 0.792 [0.781--0.803] \\
        clf$\to$filt [robust8]       & 2234 & 253 & 1653 & 0.898 & 0.575 & 0.701 [0.686--0.717] \\
        clf$\to$filt [sa32]          & 2265 & 361 & 1866 & 0.863 & 0.548 & 0.670 [0.655--0.686] \\
        \textbf{filt$\to$clf [robust8-nr]} & 2580 & 389 & 614 & \textbf{0.869} & \textbf{0.808} & \textbf{0.837 [0.827--0.849]} \\
        filt$\to$clf [robust8]       & 2182 & 243 & 1564 & 0.900 & 0.583 & 0.707 [0.691--0.725] \\
        \bottomrule
    \end{tabular}
\end{table}
% [source: thesis_eval/results_dut/tier1_results.json (dut_antiuav_960 A_bare + B_pipeline); run=thesis_eval/pipeline_eval_unified.py; ledger=dut-resolution-pipeline]

Two readings. First, dropping the reject class is what makes the fused stack competitive on a clean RGB surface: the shipped \texttt{filt$\to$clf[robust8-nr]} reaches F1 $0.837$ (recall $0.808$), recovering the recall the reject-class \texttt{robust8} surrenders here (its \texttt{filt$\to$clf} is $0.707$, recall $0.583$, because it rejects $24\%$ of these confuser-free frames). The gain is the router's; the per-frame filters remain a small recall tax on confuser-free footage, where they can only delete true positives (\texttt{filt only (mlp)} is the weakest cell at $0.726$). Second, the strongest single configuration is still RGB-alone: \texttt{ft4} reaches $F1\,0.899$ against its own ground truth (per-modality bare), above the fused $0.837$, because the cross-domain grayscale-IR channel (standalone $F1\,0.596$ in its fallback mode, far below its native thermal $F1\,0.940$ of Table~\ref{tab:rq3}) drags the fused number down. On an RGB-strong surface the fused stack buys a live second channel and a bounded precision floor rather than peak F1; the no-reject router is what keeps that cost small.

\paragraph{Why the shipped router drops the reject class.}
The four-class \texttt{robust8} can output \texttt{reject} (trust neither modality), which on a clean drone surface discards recoverable recall: it is the DUT gap above, where the router rejects $24\%$ of confuser-free frames, and the single-modality collapse of Table~\ref{tab:ablation_solo}. The shipped router removes the reject class entirely (\texttt{robust8-nr}: the same eight features, the neither-correct rows dropped in training, an argmax route to RGB, IR, or both), delegating all false-positive removal to the per-frame filter. The effect is exactly what the per-surface diagnosis predicts. On the IR-dominant Svanstr\"om and the saturated Anti-UAV, which rarely reject, it barely moves ($-1.9$~pp and $0$~pp F1 against \texttt{robust8}); on the clean RGB-strong DUT it recovers drone recall $0.575 \to 0.728$ (F1 $0.701 \to 0.792$) and on the single-modality surfaces it turns the routing collapse back into the bare detector (Table~\ref{tab:ablation_solo}). Averaged across drone surfaces the no-reject router leads on composed F1 ($0.834$ vs $0.733$). The cost is paid in confuser fire that the filter no longer has the reject net behind it to compound: RGB confuser frame fire rises $0.0015 \to 0.011$ (still $\sim$1~\%, and damped further by temporal voting, Section~\ref{sec:temporal_results}), thermal $0.217 \to 0.237$, and real-video-confuser fire $0.073 \to 0.225$ (Table~\ref{tab:temporal_production}). We judge that trade worth the recall; \texttt{robust8} is retained throughout as the ablation that quantifies what the reject class buys.

\subsection{Out-of-Distribution Confusers: the Benefit Side}
\label{sec:pipeline_confusers}

On confuser-only surfaces no drone is present; every surviving detection is a false positive, and the headline is the frame fire rate. These surfaces are evaluated in the single-modality deployment regime they represent (an RGB-only camera for the RGB confusers; a thermal camera for the IR confusers), with the router gating the one live channel.

\begin{table}[htbp]
    \centering
    \caption{Confuser-surface FP reduction. RGB confusers: $n=2{,}633$ (full test split) @640. IR confusers: $n=4{,}000$ of 5{,}237 (even spread) @640. Grayscale confusers: the same RGB confuser frames through the grayscale-fed IR channel (router bypassed; the grayscale finding's configuration, Section~\ref{sec:grayscale}). Fire $=$ fraction of frames with $\geq$1 surviving detection. The shipped \texttt{filt$\to$clf[robust8-nr]} equals \texttt{filt only (mlp)} (and equals \texttt{clf$\to$filt[robust8-nr]} to the digit, since a no-reject router never vetoes a confuser frame, so composition order is immaterial here); the \texttt{robust8} ablation reaches $0.0015$ (RGB) by adding the reject net, the suppression the shipped stack trades for recall.}
    \label{tab:ablation_confusers}
    \begin{tabular}{lccc}
        \toprule
        cell & RGB confusers & IR confusers & grayscale confusers \\
        \midrule
        bare                      & 0.304 [0.287--0.321] & 0.294 [0.280--0.309] & 0.238 [0.222--0.255] \\
        clf only [robust8-nr]     & 0.304 [0.287--0.321] & 0.294 [0.280--0.309] & --- \\
        clf only [robust8]        & 0.049 [0.041--0.057] & 0.273 [0.259--0.287] & --- \\
        clf only [robust6]        & 0.035 [0.028--0.043] & 0.223 [0.209--0.235] & --- \\
        clf only [sa32]           & 0.003 [0.001--0.005] & 0.263 [0.249--0.277] & --- \\
        filt only (mlp)           & 0.011 [0.007--0.015] & 0.237 [0.224--0.250] & 0.008 [0.004--0.011] \\
        filt only (patch)         & 0.102 [0.091--0.114] & 0.246 [0.233--0.260] & 0.204 [0.189--0.222] \\
        \textbf{filt$\to$clf [robust8-nr]} & \textbf{0.011 [0.007--0.015]} & \textbf{0.237 [0.224--0.250]} & --- \\
        clf$\to$filt [robust6]    & 0.0011 [0.000--0.003] & 0.179 [0.167--0.191] & --- \\
        clf$\to$filt [robust8]    & 0.0015 [0.000--0.003] & 0.217 [0.204--0.230] & --- \\
        \bottomrule
    \end{tabular}
\end{table}
% [source: thesis_eval/results/tier1_results.json (rgb_confuser, ir_confusers, gray_confuser C_confuser)]

Four observations. First, the \emph{filter} is the decisive stage on RGB confusers: the per-frame \texttt{mlp\_v5} cuts FP detections $835 \to 29$ ($-96.5\%$) where the patch CNN manages $835 \to 282$. The feature-reuse filter is not only $46$--$72\times$ faster than the CNN it replaced; it is $\sim$10$\times$ stronger on the surface that matters. Second, \emph{the reject class compounds with the filter, and the shipped stack trades part of that away}: the \texttt{robust8} ablation drives RGB confuser fire to $0.15\%$ ($4$ of $2{,}633$ frames; \texttt{robust6} and \texttt{sa32} compositions reach $0.11\%$ and $0.08\%$) by stacking reject on top of the filter, but the shipped no-reject \texttt{robust8-nr} leaves that work to the filter alone and sits at $1.1\%$ ($29$ of $2{,}633$): a trajectory from nearly one frame in three at the bare detector to roughly one in ninety, the deliberate cost of the recall recovered on the clean and single-modality surfaces. Third, the same filter weights transfer to the grayscale channel through nothing but the per-modality input scaler: $656 \to 21$ FP ($-96.8\%$) on grayscale-converted confusers, where the thermal scaler on the same frames manages only $656 \to 280$. The per-modality standardisation of Section~\ref{sec:ir_xmodal_verifier} is load-bearing, not a detail. (\texttt{sa32}'s near-zero router-only rate in the RGB column partly reflects its conservatism when the second modality is absent, the same conservatism that costs it recall on operational video, Section~\ref{sec:classifier_results}.)
% [source: thesis_eval/results/tier1_results.json (rgb_confuser, gray_confuser C_confuser)]

Fourth, the limit: \emph{thermal confusers resist}. On the \texttt{IR\_confusers} surface the best composition (the \texttt{robust6} ablation, $0.179$) cuts fire by only $39\%$ from bare, and the shipped no-reject cell ($0.237$) by only $19\%$. The surface is airplane-dominated ($76\%$ of frames) and the airplane category is every filter generation's weakest (Section~\ref{sec:verifier_results}); the aligned filter's confuser training is also grayscale-harvested rather than thermal-native, and these Roboflow-sourced thermal airplanes are OOD to it. In paired thermal deployments the router's modality arbitration carries more of the thermal-confuser burden than the filter does; closing the thermal airplane gap with native thermal confuser crops is recorded as future work (Section~\ref{sec:limits}).
% [source: thesis_eval/results/tier1_results.json (ir_confusers C_confuser)]

\begin{figure}[htbp]
    \centering
    \includegraphics[width=\textwidth]{fig_pipeline_ablation}
    \caption{Pipeline ablation at a glance, drawn for the reject-class \texttt{robust8} cascade (the strongest-suppression configuration; the shipped no-reject \texttt{robust8-nr} trades part of the confuser suppression for recall, Table~\ref{tab:ablation_confusers}). (a) Drone-positive paired surfaces: the router supplies most of the F1 lift and the filter finishes it; Anti-UAV confirms no-harm. Whiskers are 95\% bootstrap CIs. (b) Confuser surfaces: bare $\to$ patch $\to$ \texttt{mlp\_v5} $\to$ router $\to$ composed frame fire rates, with near-total suppression on RGB ($30.3\% \to 0.15\%$ composed) and grayscale ($23.8\% \to 0.8\%$) and measured resistance on thermal airplanes; the shipped no-reject stack delegates suppression to the filter alone and sits at $1.1\%$ on RGB. Generated from \texttt{tier1\_results.json} by \texttt{thesis\_eval/gen\_ablation\_figure.py}.}
    \label{fig:pipeline_ablation}
\end{figure}

\subsection{Solo In-Distribution Surfaces}
\label{sec:pipeline_solo}

The in-distribution test splits and the CCTV val set are single-modality surfaces: they receive the bare detector row (Part A) and the filter-only ablation, plus the router gating the one live channel.

\begin{table}[htbp]
    \centering
    \caption{Solo-surface ablation (single modality vs own GT). IR test split: $n=4{,}000$ of 9{,}612; RGB test split: $n=4{,}000$ of 17{,}209; SelCom: full 311. The \texttt{clf} rows show the router on a single live channel: the reject-class \texttt{robust8} collapses recall (it routes against a dead second modality), while the shipped no-reject \texttt{robust8-nr} degrades gracefully to the bare detector, which is why single-modality deployments can run the full router. The filter's in-distribution recall cost is the honest counterweight to its OOD benefit.}
    \label{tab:ablation_solo}
    \begin{tabular}{llccc}
        \toprule
        surface & cell & $P$ & $R$ & F1 [95\% CI] \\
        \midrule
        IR test split & bare (\texttt{v3b})   & 0.942 & 0.981 & 0.961 [0.955--0.967] \\
        IR test split & clf [robust8]         & 0.947 & 0.923 & 0.935 [0.928--0.942] \\
        IR test split & clf [robust8-nr]      & 0.942 & 0.981 & 0.961 [0.955--0.967] \\
        IR test split & filt (aligned, thermal) & 0.943 & 0.973 & 0.958 [0.952--0.964] \\
        IR test split & filt (patch)          & 0.949 & 0.931 & 0.940 [0.933--0.947] \\
        RGB test split & bare (\texttt{ft4})  & 0.954 & 0.899 & 0.926 [0.919--0.933] \\
        RGB test split & clf [robust8]        & 0.981 & 0.383 & 0.551 [0.532--0.570] \\
        RGB test split & clf [robust8-nr]     & 0.954 & 0.899 & 0.926 [0.919--0.933] \\
        RGB test split & filt (mlp)           & 0.976 & 0.691 & 0.809 [0.797--0.823] \\
        RGB test split & filt (patch)         & 0.957 & 0.832 & 0.890 [0.881--0.898] \\
        SelCom val & bare (\texttt{ft4})      & 0.858 & 0.451 & 0.591 [0.532--0.639] \\
        SelCom val & clf [robust8]            & 0.881 & 0.125 & 0.220 [0.161--0.276] \\
        SelCom val & clf [robust8-nr]         & 0.858 & 0.451 & 0.591 [0.531--0.639] \\
        SelCom val & filt (mlp)               & 0.950 & 0.451 & 0.612 [0.551--0.659] \\
        SelCom val & filt (patch)             & 0.858 & 0.451 & 0.591 [0.532--0.639] \\
        \bottomrule
    \end{tabular}
\end{table}
% [source: thesis_eval/results/tier1_results.json (ir_dset_final, rgb_dataset_test, selcom_val S4_verifier)]

Three per-surface readings. On the IR test split the aligned filter costs $0.8$~pp recall for a precision nudge (recall-safe, as designed), while the patch CNN costs $5.0$~pp recall on the same frames. On the RGB test split the feature-reuse filter's known carve-out reproduces exactly at the canonical configuration: $-11.7$~pp F1, all of it recall ($0.899 \to 0.691$), the OOD coverage gap diagnosed in Section~\ref{sec:verifier_results}. On SelCom the filter cuts false positives $22 \to 7$ at \emph{zero} recall cost (precision $0.858 \to 0.950$), while the patch CNN, which has never seen CCTV-scale crops, is bit-for-bit the bare detector: the abstention behaviour that motivated the feature-reuse design in the first place. A regime note: the reject-class \texttt{robust8} collapses recall on these single-modality surfaces (it routes against a dead second channel; clf-only $R=0.383$ on the RGB test split, $0.125$ on SelCom), but the shipped no-reject \texttt{robust8-nr} does not, instead degrading to the bare detector ($R=0.899$ and $0.451$, the \texttt{clf [robust8-nr]} rows above). The collapse was the reject class, not routing per se, so single-modality deployments can run the shipped router safely; detector $+$ filter remains the minimal configuration the GUI's single-camera modes use.
% [source: thesis_eval/results/tier1_results.json (S4_verifier clf rows)]

\subsection{Modality A/B: What a Single-Camera System Misses}
\label{sec:modality_ab}

RQ3 asks whether routed dual-modality beats either modality alone. The headline trust-aware rule cannot answer that question directly, because it deliberately excludes the distrusted side's ground truth (Section~\ref{sec:metrics}); a single-camera system, however, \emph{physically misses} whatever only the other sensor sees. This ablation therefore re-scores the same cached detections under a \emph{coverage} rule: every modality's ground truth always counts, and a missing or distrusted modality's targets are false negatives. All arms use the same rule, so they are directly comparable; the numbers are deliberately lower than, and not comparable to, the trust-aware headline cells.

\begin{table}[htbp]
    \centering
    \caption{Modality A/B under coverage scoring (absent/distrusted side's GT counts as missed). Svanstr\"om paired, $n=4{,}000$, IoP@0.5; Anti-UAV paired, $n=4{,}000$, IoU@0.5. ``+filt'' applies the per-modality MLP confuser filter.}
    \label{tab:modality_ab}
    \begin{tabular}{lcc}
        \toprule
        arm & Svanstr\"om F1 [95\% CI] & Anti-UAV F1 [95\% CI] \\
        \midrule
        RGB only            & 0.458 [0.448--0.468] & 0.643 [0.640--0.646] \\
        RGB only $+$filt    & 0.582 [0.573--0.590] & 0.643 [0.640--0.646] \\
        IR only             & 0.632 [0.626--0.637] & 0.652 [0.648--0.656] \\
        IR only $+$filt     & 0.632 [0.626--0.637] & 0.652 [0.648--0.656] \\
        both (unrouted)     & 0.742 [0.730--0.754] & 0.973 [0.970--0.976] \\
        both $+$filt        & 0.907 [0.900--0.915] & 0.973 [0.970--0.976] \\
        routed [\texttt{robust8}]        & 0.915 [0.908--0.921] & 0.973 [0.970--0.976] \\
        routed [\texttt{robust8}] $+$filt & 0.921 [0.914--0.927] & 0.973 [0.970--0.976] \\
        routed [\texttt{robust8-nr}]     & 0.743 [0.731--0.755] & 0.974 [0.971--0.976] \\
        \textbf{routed [\texttt{robust8-nr}] $+$filt (production)} & \textbf{0.909} [0.902--0.916] & \textbf{0.974} [0.971--0.976] \\
        \bottomrule
    \end{tabular}
\end{table}
% [source: thesis_eval/results/notes_round1_results.json (M_modality_ab); run=thesis_eval/notes_round1_replays.py]

The answer is unambiguous on the discriminating surface. On Svanstr\"om, either camera alone misses roughly half of what the pair can see (RGB-only $0.458$, IR-only $0.632$); carrying both sensors unrouted recovers the coverage but imports both failure profiles ($0.742$); and routing more than recovers it: the shipped no-reject \texttt{robust8-nr}$+$filt reaches $0.909$, more than $0.27$ above either single camera. Its margin over both$+$filt is small here ($+0.2$~pp, $0.909$ vs $0.907$), because without a reject class it rarely overrides the both-channel default on this surface; the reject-class \texttt{robust8} ablation widens that margin to $+1.4$~pp ($0.921$) by vetoing frames the filter cannot, the coverage cost it trades for the recall recovered elsewhere. On the in-distribution Anti-UAV control every dual arm is statistically identical ($0.973$), the expected no-harm result; the single-modality arms sit near $0.65$ there only because each sensor sees just its own half of the paired ground truth. The asymmetric filter effect is also visible: the filter moves RGB-only by $+12.4$~pp (it removes the visible channel's confuser false positives) and IR-only by nothing (the thermal channel barely hallucinates on this surface, Section~\ref{sec:failure_profile}).

\subsection{Per-Size Attribution: Where the Resolvable Floor Bites}
\label{sec:per_size}

Several claims in this thesis lean on drone \emph{size}: the resolution floor of Section~\ref{sec:resolution_arch}, the \texttt{retrained\_v2} small-drone collapse, and the confuser filter's coverage carve-out. The unified cache stores every ground-truth and detection box, so per-size precision/recall replays exactly; buckets are ground-truth $\sqrt{\text{area}}$ in native-image pixels, with TP/FN bucketed by the ground-truth box and FP by the predicted box. A bucket's recall is only reported when it contains ground truth.

\begin{table}[htbp]
    \centering
    \caption{Per-size recall, bare $\to$ $+$filter, on the two surfaces that carry size-dependent claims (production detectors, Tier-1 configurations). The RGB-test carve-out is concentrated below 32~px; Svanstr\"om IR ground truth is dominated by sub-16~px drones, which is why \texttt{imgsz} matters there.}
    \label{tab:per_size}
    \begin{tabular}{llcccc}
        \toprule
        surface & detector & bucket & $n_{gt}$ & bare $R$ & $+$filt $R$ \\
        \midrule
        rgb\_dataset test & \texttt{ft4}  & $<$16 px  & 726   & 0.782 & 0.256 \\
                          &               & 16--32 px & 607   & 0.865 & 0.447 \\
                          &               & 32--64 px & 702   & 0.936 & 0.829 \\
                          &               & $\geq$64 px & 1{,}417 & 0.956 & 0.951 \\
        \midrule
        Svanstr\"om       & \texttt{ft4} (RGB)  & $<$16 px  & 27    & 0.630 & 0.593 \\
                          &               & 16--32 px & 994   & 0.897 & 0.823 \\
                          &               & 32--64 px & 633   & 0.956 & 0.918 \\
        \midrule
        Svanstr\"om       & \texttt{v3b} (IR)   & $<$16 px  & 993   & 0.970 & 0.970 \\
                          &               & 16--32 px & 643   & 0.998 & 0.998 \\
        \bottomrule
    \end{tabular}
\end{table}
% [source: thesis_eval/results/notes_round1_results.json (SZ_per_size); run=thesis_eval/notes_round1_replays.py]

Two attributions emerge. First, the confuser filter's RGB-test regression (Section~\ref{sec:mlp_recall_drop}) is now localised: it is a \emph{small-drone} recall cost ($0.782 \to 0.256$ below 16~px, $0.865 \to 0.447$ at 16--32~px) while drones above 64~px are untouched ($0.956 \to 0.951$). The filter's training corpus under-covers tiny crops, and tiny crops are exactly where detector features are weakest. Second, the resolvable-floor story is quantified: Svanstr\"om's median drone is $29.8$~px ($\sqrt{\text{area}}$, native $640{\times}512$ RGB) and the IR ground truth is majority sub-16~px (median $14.8$~px), so a $640$-pixel inference side leaves the median IR drone $\sim$15 native pixels across; this is the measured basis for the \texttt{imgsz=1280} requirement on the Svanstr\"om RGB channel (Section~\ref{sec:canonical_configs}).

\subsection{Background Failure Profile}
\label{sec:failure_profile}

Where do the remaining errors live: in the scene, or on the object? All 273 Svanstr\"om source sequences and the 9 drone video clips were manually tagged with a coarse background class from one representative frame each (\texttt{sky}: near-pure sky; \texttt{horizon}: sky with a visible ground/treeline/building strip; \texttt{ground}: clutter-dominant); the tag propagates to every cached frame of the sequence, and per-background scores replay from the cache. The tags are coarse, single-frame, and single-annotator, and sequences are split into drone-bearing and confuser-only before aggregating, because Svanstr\"om's pure-sky sequences are disproportionately its helicopter and airplane clips, a confound that would otherwise dominate the table.

\begin{table}[htbp]
    \centering
    \caption{Background failure profile on Svanstr\"om (bare detectors, per modality vs own GT, Tier-1 rules). On confuser sequences the fire rate barely moves with background; on drone sequences precision is high on every background and recall dips on pure sky.}
    \label{tab:failure_profile}
    \begin{tabular}{llcccc}
        \toprule
        sequences & background & $n$ frames & RGB FP-frame rate & RGB $R$ & IR FP-frame rate \\
        \midrule
        confuser-only & sky     & 1{,}294 & 0.667 & --- & 0.006 \\
        confuser-only & horizon & 888     & 0.741 & --- & 0.017 \\
        confuser-only & ground  & 147     & 0.646 & --- & 0.000 \\
        \midrule
        drone-bearing & sky     & 394     & 0.074 & 0.827 & 0.102 \\
        drone-bearing & horizon & 1{,}114 & 0.033 & 0.944 & 0.073 \\
        drone-bearing & ground  & 163     & 0.000 & 0.933 & 0.153 \\
        \bottomrule
    \end{tabular}
\end{table}
% [source: thesis_eval/results/failure_profile_results.json; run=thesis_eval/failure_profile_aggregate.py (282 manual seq tags, 2026-06-11)]

The profile supports the thesis framing directly. On confuser-only sequences the RGB detector fires on $65$--$74\%$ of frames \emph{regardless of background}: the hallucination is driven by the confuser object, not by scene clutter, which is why a per-detection filter (which looks at the object) works and scene-level heuristics would not. On drone-bearing sequences the same detector is clean ($0$--$7\%$ FP-frame rate) on every background, and its recall dips on pure-sky sequences ($0.827$ vs $0.944$ on horizon), consistent with the per-size table: the pure-sky clips carry the most distant drones. The thermal channel inverts the profile: near-zero fire on confuser sequences ($\leq 1.7\%$, airplane-dominated) and a mild ground-clutter penalty on drone sequences ($15.3\%$ FP-frame rate), which is one of the complementarities the trust router exploits. On the video clips the same breakdown is confounded with per-clip drone distance and is reported only as context in the results file.

\subsection{The Filter Buys Back the Detector Floor}
\label{sec:lowconf_mode}

The production confidence floors (\texttt{rgb\_conf}$=$0.25, \texttt{ir\_conf}$=$0.40) exist for one reason: to keep the bare detector's false positives in check. With false-positive control delegated to the per-frame filter, that reason weakens. A detector-confidence sweep ($0.05$--$0.50$, crossed with the filter, replayed from a dedicated floor-$0.05$ cache) therefore asks whether \emph{lowering the floor and letting the filter absorb the flood} beats the defaults.

On the deployment-anchor surface it does. On SelCom (full $311$-image val), dropping the floor to $0.05$ raises recall $0.451 \to 0.678$ ($+22.7$~pp) and the filter restores the precision the low floor surrenders ($0.531 \to 0.707$), for $F1 = 0.692$ $[0.654$--$0.730]$ against the default's bare $0.591$ and filtered $0.612$ $[0.551$--$0.659]$. The filter is \emph{recall-transparent} here (filtered recall is identical to bare recall at every floor), which is precisely the mechanism the design promises: the detector proposes, the filter disposes, and nothing true is lost. The move is also confuser-safe: at floor $0.05$ the RGB-confuser corpus floods the bare detector with $1{,}281$ FP detections and the filter holds the line at $45$ ($96.5\%$ suppression; the default operating point leaves $29$). One bound applies: the gain over the \emph{filtered} default is marginal at the CI level (the intervals overlap), while the gain over the bare default is clean.
% [source: ledger=recall-plus-filter-ood-mode; eval=selcom_confsweep_lowconf_mlp,rgb_confuser_confsweep_lowconf; run=thesis_eval/conf_sweep_replay.py; cache=thesis_eval/results/conf_sweep/conf_sweep_results.json]

\begin{table}[htbp]
    \centering
    \caption{Detector-floor sweep $\times$ filter on SelCom (full 311 images, IoP@0.5, \texttt{imgsz=1280}). The filter is recall-transparent at every floor; lowering the floor converts directly into recall.}
    \label{tab:lowconf_selcom}
    \begin{tabular}{lcc}
        \toprule
        detector floor & bare $P/R/F1$ & $+$\texttt{mlp\_v5} $P/R/F1$ \\
        \midrule
        0.25 (default) & 0.858 / 0.451 / 0.591 & 0.950 / 0.451 / 0.612 \\
        0.10           & 0.658 / 0.573 / 0.612 & 0.816 / 0.573 / 0.673 \\
        \textbf{0.05}  & 0.531 / 0.678 / 0.595 & \textbf{0.707 / 0.678 / 0.692} \\
        \bottomrule
    \end{tabular}
\end{table}
% [source: eval=selcom_confsweep_lowconf_mlp; run=thesis_eval/conf_sweep_replay.py]

The mode does \emph{not} generalise to a new global default, and the sweep maps its boundary precisely. On in-domain RGB the filter's coverage carve-out dominates at every floor ($-12.5$~pp vs the bare default; keep bare at $0.25$ there). On the thermal side the floor must stay at $0.40$: dropping it balloons bare thermal-confuser FPs $1{,}051 \to 1{,}911$ and the aligned filter absorbs only $\sim$18--20\% of thermal confusers (Table~\ref{tab:ablation_confusers}). On the saturated Anti-UAV pipeline (a $500$-frame screening subsample) every cell sits within CI of the default's $0.976$; there is nothing to gain. The result is therefore an \emph{operating mode}, not a setting: on hard out-of-distribution RGB surfaces where the detector is recall-starved, drop the RGB floor to $0.05$--$0.10$ and let the filter own precision; everywhere else the defaults stand. One regime caveat travels with it: the trust router's features were mined at the $0.25$ floor, so sub-floor router cells are deployed-behaviour estimates rather than retraining claims.
% [source: ledger=recall-plus-filter-ood-mode; eval=antiuav_confsweep_pipeline,ir_confsweep_lowconf; run=thesis_eval/conf_sweep_replay.py; cache=docs/analysis/2026-06-11_conf_sweep_recall_plus_filter.md]

\subsection{Runtime Cost}
\label{sec:pipeline_speed}

\begin{table}[htbp]
    \centering
    \caption{Per-stage runtime (GTX 1050 Ti). The production stages are cheap enough to run on every frame; the predecessors they replaced were the reason gating architectures existed at all.}
    \label{tab:speed}
    \begin{tabular}{llll}
        \toprule
        stage & predecessor & production & speedup \\
        \midrule
        trust router    & \texttt{fusion\_no\_fn} 38.3 ms/frame & \texttt{robust8-nr} 0.095 ms/frame & $\sim$404$\times$ \\
        confuser filter & patch CNN 59--112 ms/det & \texttt{mlp\_v5} 1.3--2.1 ms/det & 37--72$\times$ \\
        \bottomrule
    \end{tabular}
\end{table}
% [source: ledger=robust6-speed-feature-efficiency,v5-ship-per-frame; run=eval/bench_speed.py; cache=docs/analysis/2026-06-04_consolidated_sell_new_stack.md]

Total pipeline overhead beyond the two detector forward passes is $1$--$4\%$. This is what dissolves the alert-gating compromise of the predecessor design (Section~\ref{sec:alert_gate}): when the filter costs two milliseconds, it simply runs always.

% ======================================================================
\section{Temporal Evidence at the Segment Grain}
\label{sec:temporal_results}

The deployed system does not alert per frame: an $N$-of-$M$ smoother (Section~\ref{sec:temporal}) aggregates per-frame decisions into alert events. The segment-level evaluation runs on the 19-clip real-video surface, the one canonical surface with consecutive frames ($1{,}359$ drone-clip frames across 9 clips; $1{,}250$ confuser-clip frames across 10), using sliding 3-frame windows with a 2-of-3 vote per clip, in the dual-channel video regime (\texttt{ft4} on the RGB frame, \texttt{v3b} on its grayscale conversion, \texttt{is\_grayscale}$=$1).

\begin{table}[htbp]
    \centering
    \caption{Segment-level (2-of-3 windows) evaluation on the real-video clips. Fully out-of-distribution for every component. The per-frame filter thresholds are the benchmark-tuned production values; the last column is the relative confuser-fire cut from bare. The patch arm here is a per-frame approximation of the alert-gated original.}
    \label{tab:temporal_production}
    \begin{tabular}{lcccc}
        \toprule
        cell & window $P$ & window $R$ & window F1 & confuser window fire (cut) \\
        \midrule
        bare (both channels)            & 0.957 & 0.753 & 0.843 & 0.350 (---) \\
        filt only (mlp)                 & 0.954 & 0.541 & 0.690 & 0.226 ($-35\%$) \\
        filt only (patch, per-frame)    & 0.957 & 0.689 & 0.801 & 0.182 ($-48\%$) \\
        clf only [robust8-nr]           & 0.957 & 0.753 & 0.843 & 0.350 ($-0\%$) \\
        clf only [robust8]              & 0.967 & 0.520 & 0.676 & 0.101 ($-71\%$) \\
        clf only [robust6]              & 0.999 & 0.584 & 0.737 & 0.081 ($-77\%$) \\
        clf only [sa32]                 & 0.967 & 0.582 & 0.727 & 0.120 ($-66\%$) \\
        clf$\to$filt [robust8-nr]       & 0.954 & 0.537 & 0.687 & 0.225 ($-36\%$) \\
        clf$\to$filt [robust8]          & 0.972 & 0.393 & 0.560 & 0.073 ($-79\%$) \\
        clf$\to$filt [robust6]          & 0.998 & 0.443 & 0.614 & 0.057 ($-84\%$) \\
        clf$\to$filt patch [sa32]       & 0.969 & 0.535 & 0.689 & 0.080 ($-77\%$) \\
        \bottomrule
    \end{tabular}
\end{table}
% [source: thesis_eval/results/temporal_results.json; run=thesis_eval/temporal_replay.py (window 2-of-3, patch_thr=0.7, mlp rgb@0.25/gray@0.25)]
% [nr rows: thesis_eval/results_noreject/temporal_results.json — clf[robust8-nr]=bare; clf->filt[robust8-nr]=filt-only (no-reject router passes every confuser frame to the filter)]

This surface is fully OOD for every component, and it is where the pipeline's costs show. Three results. \emph{First, the trade is real on operational video, and the shipped no-reject router sits on the recall side of it:} \texttt{clf$\to$filt[robust8-nr]} keeps window recall $0.537$ (F1 $0.687$, the best composed F1 here) while cutting confuser fire $36\%$, whereas the reject-class compositions (\texttt{robust8}/\texttt{robust6}) push the cut to $79$--$84\%$ but veto OOD drones aggressively (window recall $0.39$--$0.44$), the video counterpart of the filter's known coverage gap (Section~\ref{sec:verifier_results}). The per-frame filter thresholds were tuned on benchmark surfaces, so on fully OOD video they over-veto, which is exactly where the deployment's alert-gate placement matters. \emph{Second, the gate position matters more than the gate:} the predecessor design consulted its filter only at the alert moment, and under that placement the same surface showed segment F1 \emph{gains} (below); a per-frame veto applies the same discrimination thousands of times more often, so every false veto compounds. The measured per-frame numbers therefore \emph{validate the GUI's alert-gate placement} for OOD video rather than contradicting the pipeline: routing per-frame (cheap, $-71$--$77\%$ fire at $10$--$17$~pp F1) plus verification at the alert gate is the operating point the deployment uses. \emph{Third, temporal voting is recall-neutral and precision-positive here} (window vs frame: $\Delta R \approx \pm 0.01$ on every cell, confuser fire consistently lower), confirming the smoother's role as a stabiliser rather than a recall recoverer when detections are already temporally consistent.

A filter-threshold sweep over the cached probabilities settles whether the veto cost is an operating-point artifact: it is not. Lowering the shared filter threshold from $0.25$ to $0.01$ recovers most of the composed cells' recall \emph{relative to router-only} (robust8: $R$ $0.393 \to 0.507$ against the router's $0.520$; robust6: $0.443 \to 0.577$ against $0.584$): at a near-zero threshold the per-frame filter costs only $\sim$1~pp recall on top of the router while still trimming window fire by $7$--$9\%$ relative. But the true drones' filter probabilities on this surface are smeared across the $[0.01, 0.25)$ interval rather than concentrated near 1, so no threshold simultaneously keeps them and rejects confusers: the video-deployment choice is a near-zero per-frame threshold (cheap, marginal benefit) or the alert gate (the predecessor's measured-gain placement). The full sweep is \texttt{thesis\_eval/results/video\_thr\_sweep.md}.
% [source: thesis_eval/results/video_thr_sweep.json; run=thesis_eval/video_thr_sweep.py]

\paragraph{Design-evolution evidence (predecessor stack).}
The temporal smoother's value was first established on this surface with the predecessor stack (baseline RGB, \texttt{sa32}, alert-gated patch filter); those numbers are retained as the design-evolution record. At the segment grain that cascade \emph{gained} drone $F1$ for all four RGB variants tested ($+6.6$ to $+16.5$~pp; baseline RGB $0.760 \to 0.826$) while cutting confuser FPR $39$--$81\%$ (baseline: $0.512 \to 0.162$), and it compressed the spread between detector variants from $15.5$ to $5.6$~pp F1: the cascade absorbs detector variance. Per-frame metrics on the same runs \emph{fell} for three of four variants, which is the wrong-unit-of-analysis lesson: the per-frame grain counts every wrong box, the deployed system counts alerts.
% [source: ledger=cascade-segment-recovers,cascade-tightens-variance,cascade-perframe-misleading; eval=pipe_vid_baseline_seg,pipe_vid_baseline_pf; config=pipe_video_drone_iop (predecessor configuration: baseline RGB + sa32 + alert-gated patch @0.7)]

\begin{figure}[htbp]
    \centering
    \includegraphics[width=\textwidth]{fig8_cascade_segment}
    \caption{Design-evolution record (predecessor stack: baseline RGB, \texttt{sa32}, alert-gated patch): segment-level drone $F1$ and confuser FPR, Stage-1 RGB (grey) versus the full cascade, on the real-video diagnostic. (a) The cascade lifts drone $F1$ for every variant. (b) It cuts confuser FPR by $39$--$81\%$. The production-stack re-run of this surface is Table~\ref{tab:temporal_production}.}
    \label{fig:cascade_segment_fig}
\end{figure}

The per-category breakdown of the predecessor runs carries one durable lesson: the cascade's suppression concentrates on birds ($45$--$57\times$ on the worst variants), is moderate on helicopters, and is essentially inert on airplanes ($\sim$0.21 FPR across all variants); the patch filter is genuinely uncertain on airplane crops (median confuser probability $0.540$ vs $0.904$ for birds). Airplanes remain the weakest category for every filter generation (Section~\ref{sec:verifier_results}).
% [source: ledger=cascade-bird-vs-airplane-asymmetry; eval=pipe_percat_sa32; config=pipe_video_confuser]

% ======================================================================
\section{Component Deep-Dives}
\label{sec:components_eval}

\subsection{RGB Detector}
\label{sec:rgb_results}

\begin{table}[htbp]
    \centering
    \caption{RGB variants on Svanstr\"om @ \texttt{imgsz=1280} (IoP@0.5, drone class; full 28{,}710 frames, the pre-protocol full-set run, internally consistent). Aggressive retraining destroys recall on small Svanstr\"om drones while leaving its own in-domain test performance intact.}
    \label{tab:rgb_comparison}
    \begin{tabular}{lccccc}
        \toprule
        Model & Drone $P$ & Drone $R$ & Bird halluc & Airplane halluc & Heli halluc \\
        \midrule
        baseline           & 0.940 & \textbf{0.961} & 94.4\% & 74.6\% & 66.2\% \\
        \texttt{hardneg\_v3more}  & 0.941 & 0.950 & 94.2\% & 64.7\% & 41.9\% \\
        \texttt{retrained\_v2}    & 0.943 & 0.306 & \textbf{3.4\%}  & \textbf{5.6\%}  & \textbf{4.5\%}  \\
        \bottomrule
    \end{tabular}
\end{table}
% [source: ledger=retrainedv2-recall-collapse; eval=rgb_svan_baseline,rgb_svan_hardneg,rgb_svan_retrainedv2; run=eval/diagnose_failures_all.py; cache=eval/results/_failure_diagnosis/svanstrom_1280_by_category.csv; config=svan_iop_1280 (full set)]

The three stances trace the lever that motivated the whole architecture (Section~\ref{sec:rgb_three_variants}): hard-negative mining is class-specific (helicopters yield, birds do not), and the one stance that suppresses birds collapses drone recall. The dedicated resolution sweep (Section~\ref{sec:resolution_arch}, Table~\ref{tab:resolution}) puts the \texttt{imgsz=640} regime in within-model terms: the baseline loses $28$~pp of recall at 640 and \texttt{retrained\_v2} drops further from its already-collapsed level. (That sweep reports \texttt{retrained\_v2} at $R=0.323$ at 1280 against the $0.306$ here: the same already-collapsed level measured on its stride-7 harness, $n=4{,}102$, rather than this full $28{,}710$-frame set, a sampling difference of $1.7$~pp, not an inconsistency.)
% [source: eval/results/svan_resolution_sweep.json; run=eval/svan_resolution_sweep.py] The collapse does not show in-domain: on the composite corpus's \emph{own} test split the same weights score $F1=0.949$, above the baseline's $0.942$, and on SelCom they total-collapse to $F1=0.007$ (1~TP of 295~GT). In-domain evaluation alone would have selected the detector that fails deployed; this is the sharpest single argument for the OOD surfaces the protocol mandates.
% [source: ledger=retrainedv2-recall-collapse,rgb-collapse-ood-specific; eval=rgb_rgbds_baseline,rgb_rgbds_retrainedv2,selcom_retrainedv2_1280; config=rgb_dataset_iou_640,selcom_iop_1280]

\begin{table}[htbp]
    \centering
    \caption{SelCom CCTV fine-tune; IoP@0.5, conf=0.25, val split 311~imgs / 295~GT. Doubling \texttt{imgsz} doubles recall \emph{and} raises precision: a resolution win, not a trade.}
    \label{tab:selcom}
    \begin{tabular}{lcccccc}
        \toprule
        Model & imgsz & TP & FP & FN & $P$ & $R$ \\
        \midrule
        baseline (\texttt{Yolo26n\_trained})        & 640  &   2 &  6 & 293 & 0.250 & 0.007 \\
        baseline                                    & 1280 &  26 & 37 & 269 & 0.413 & 0.088 \\
        \texttt{Yolo26n\_selcom\_mixed\_ft2}        & 640  &  72 & 50 & 223 & 0.590 & 0.244 \\
        \textbf{\texttt{Yolo26n\_selcom\_mixed\_ft2\_1280}} & \textbf{1280} & \textbf{138} & \textbf{43} & \textbf{157} & \textbf{0.762} & \textbf{0.468} \\
        \bottomrule
    \end{tabular}
\end{table}
% [source: ledger=selcom-imgsz-win; eval=selcom_ft2_1280; run=RGB model/finetune_selcom.py; cache=runs/rgb_finetune_eval/Yolo26n_selcom_mixed_ft2_1280/comparison.json; config=selcom_iop_1280]

At \texttt{imgsz=960}, \texttt{selcom\_960} is the best SelCom-holdout drone-first RGB variant tested ($F1=0.585$ on the hold-out); it is deliberately not in the production stack (it loses to the baseline on Svanstr\"om medium-range drones) but is recorded as the deployment operating point for this camera.
% [source: ledger=selcom960-cross-surface-winner; eval=xsurf_selcom_holdout; config=selcom_holdout_persize]

\paragraph{Real-video cameo: the three-stance continuum generalises.}
On the 19-clip diagnostic (per-clip table in Appendix~\ref{app:datasets}), the three stances occupy the same continuum (\texttt{retrained\_v2} high-precision/low-recall at $P=0.909, R=0.453$; \texttt{selcom\_1280} high-recall/low-precision at $P=0.649, R=0.812$; baseline between), and no single RGB variant dominates across scenes. The full six-mode table (including the IR-on-grayscale rows it introduced) is Table~\ref{tab:realvideo_master} in Section~\ref{sec:grayscale}.
% [source: ledger=no-pareto-dominance-video; eval=vid_drone_baseline; config=video_drone_iop]

\subsection{IR Detector: the Six-Revision HITL Case Study}
\label{sec:ir_evolution}

The IR detector was developed across five numbered revisions (V2 through V6) plus a \texttt{Final} retrain and a \texttt{v3b} corrective fine-tune, by the loop of Section~\ref{sec:hitl_method}. Earlier per-version best-epoch mAP50 numbers are not directly comparable because the val split changed with the dataset. Every version was therefore re-evaluated on the same fixed \texttt{IR\_dset\_final} test split at \texttt{imgsz=640} (script: \texttt{eval/eval\_ir\_versions.py}).

\begin{table}[htbp]
    \centering
    \caption{IR detector versions evaluated on the fixed \texttt{IR\_dset\_final} test split @ \texttt{imgsz=640}. Comparable across versions; supersedes the per-version best-epoch mAP narrative. \texttt{v3b} is the production weight.}
    \label{tab:ir_evolution}
    \begin{tabular}{lcccccc}
        \toprule
        Version & mAP50 & mAP50-95 & $P$ & $R$ & F1 & Status \\
        \midrule
        V2     & 0.458 & 0.219 & 0.661 & 0.406 & 0.503 & superseded \\
        V3     & 0.571 & 0.242 & 0.648 & 0.579 & 0.611 & superseded \\
        V4     & 0.722 & 0.393 & 0.895 & 0.669 & 0.765 & superseded \\
        V5     & 0.694 & 0.446 & 0.768 & 0.709 & 0.737 & superseded (regression) \\
        V6     & 0.956 & 0.571 & 0.921 & 0.941 & 0.931 & superseded \\
        Final  & \textbf{0.977} & \textbf{0.602} & 0.955 & \textbf{0.980} & 0.967 & base for prod \\
        \textbf{v3b (production)} & 0.972 & 0.591 & \textbf{0.957} & 0.977 & \textbf{0.967} & current \\
        \bottomrule
    \end{tabular}
\end{table}

\begin{figure}[htbp]
    \centering
    \includegraphics[width=0.85\textwidth]{fig4_ir_evolution}
    \caption{IR detector P/R trajectory across the five numbered revisions (V2--V6) plus the \texttt{Final} retrain and production \texttt{v3b}, on the fixed \texttt{IR\_dset\_final} test split. The V5 regression (precision $0.895 \to 0.768$) is recovered by V6 and the production \texttt{v3b}.}
    % [source: ledger=ir-version-progression; eval=ir_final_v3b; run=eval/ir_version_comparison.py; cache=eval/results/ir_version_comparison/ir_comparison_test_640_2026-05-16T20-04-39.csv; config=ir_final_640]
    \label{fig:ir_evolution}
\end{figure}

The trajectory reads in four passes. \emph{V2 $\to$ V3}: a narrower curated subset trades a little precision for a large recall gain ($0.406 \to 0.579$), the premise for the mining cycles that follow. \emph{V3 $\to$ V4}: precision jumps $0.648 \to 0.895$ with recall barely moving, the signature of an FP-mining cycle. \emph{V5, the regression}: precision drops to $0.768$ because a large positive batch was injected all at once, bypassing the review loop. The un-reviewed batch carried unlabelled drones; V5 fires on them and is scored as FP. The model did not get worse; the dataset state did, which is precisely the coupling Section~\ref{sec:coevolution} warns about. \emph{V6/\texttt{Final}}: corpus cleanup recovers and surpasses ($F1=0.967$). \texttt{Final} and \texttt{v3b} are indistinguishable on this split; \texttt{v3b} is retained as production because its 2-epoch corrective fine-tune addresses deployment failures (GUI regressions, hot-pixel scenes) not represented here.
% [source: ledger=ir-version-progression; eval=ir_final_v2..ir_final_v3b; run=eval/ir_version_comparison.py; cache=eval/results/ir_version_comparison/ir_comparison_test_640_2026-05-16T20-04-39.csv; config=ir_final_640]

\paragraph{OOD bound.} The IR detector's strength is in-distribution. Its off-distribution \emph{recall} on independent thermal drone footage is untested on verified data (exploratory runs on community thermal sets suggested substantial degradation, but those sets' label quality could not be verified and their numbers are not reported); this is a stated limitation (Section~\ref{sec:threats}), and one of the reasons the RGB detector is retained as a parallel channel rather than treating thermal as primary.

\subsection{Trust Classifier}
\label{sec:classifier_results}
\label{sec:classifier_compare}

The production router is the no-reject \texttt{robust8-nr}; this section reports the design path that produced it, from statistical feature selection through the reject-class variants to the shipped no-reject form. The lineage is one feature set with one decision rule changed. Feature selection (Section~\ref{sec:feature_selection}) produced the six-feature \texttt{robust6} and its eight-feature extension \texttt{robust8}, both \emph{four-class} routers that can label a frame \texttt{reject} (trust neither modality). The shipped router keeps those eight features but \emph{drops the reject class}, delegating all false-positive removal to the per-frame filter. The reject-class variants flank each other in the pipeline ablation (Tables~\ref{tab:ablation_svanstrom}--\ref{tab:ablation_confusers}) exactly as the feature statistics predicted (\texttt{sa32} marginally ahead in-domain, the statistical generation ahead out-of-domain and $404\times$ cheaper); the shipped \texttt{robust8-nr} then trades a little of \texttt{robust8}'s confuser suppression for the recall it recovers on clean, single-modality, and degraded-channel surfaces.

\paragraph{The statistical generation, validated end-to-end.}
\texttt{robust6} (six leakage-free features, Section~\ref{sec:feature_selection}) was validated against \texttt{sa32} in a 5{,}000-frame/surface alert ablation under the production cascade: in-domain Svanstr\"om F1 within $0.2$~pp of \texttt{sa32} ($0.9957$ vs $0.9974$, classifier$\to$filter cells), and materially better on the OOD confuser surface (classifier-only fire $0.143$ vs $0.203$, $\approx$30\% fewer false alerts; best-composed cascade $0.057$ vs $0.066$). The same run makes a structural point: on surveillance-style footage the classifier does the bulk of FP suppression (classifier-only fire $0.003$ vs $0.065$ for the filter alone), while on confuser-rich footage the two stages are complementary.
% [source: ledger=robust6-production-viable; run=eval/overnight_ablation.py; cache=eval/results/_overnight_ablation/ablation_results.json; config=clf_own_holdout (5k strided frames/surface)]

\begin{table}[htbp]
    \centering
    \caption{\texttt{robust6} vs \texttt{sa32} under the production cascade ($5{,}000$ strided frames/surface). \texttt{sa32} edges in-domain; \texttt{robust6} wins out-of-domain at $1/5$th the feature count and none of the OpenCV cost.}
    \label{tab:robust6_pipeline}
    \begin{tabular}{lcc}
        \toprule
        cell & Svanstr\"om $F1$ $\uparrow$ & confuser fire $\downarrow$ \\
        \midrule
        bare (no classifier/filter)        & 0.715 & 0.379 \\
        filter only                        & 0.964 & 0.114 \\
        classifier only \texttt{[sa32]}      & 0.9972 & 0.203 \\
        classifier only \texttt{[robust6]}   & 0.9877 & \textbf{0.143} \\
        classifier$\to$filter \texttt{[sa32]}    & \textbf{0.9974} & 0.079 \\
        classifier$\to$filter \texttt{[robust6]} & 0.9957 & 0.066 \\
        filter$\to$classifier \texttt{[robust6]} & --- & \textbf{0.057} \\
        \bottomrule
    \end{tabular}
\end{table}

\paragraph{The shipped router: dropping the reject class.}
The reject class is a liability on any surface that is not a confuser-rich paired stream. When a modality is absent (single-camera deployment), degraded (the grayscale-IR fallback), or the surface is simply clean (DUT), \texttt{reject} fires on recoverable drone frames and recall collapses: clf-only $R\,0.383$ on the RGB test split, $0.125$ on SelCom, and a $24\%$ reject rate on confuser-free DUT (Tables~\ref{tab:ablation_solo},~\ref{tab:ablation_dut}). \texttt{robust8-nr} removes the class at its source: it retrains the eight-feature \texttt{robust8} with the neither-correct rows dropped, so the router always routes to RGB, IR, or both (argmax, no threshold) and the per-frame filter owns all false-positive rejection. This recovers the lost recall in full (composed F1: DUT $0.701\to0.792$, RGB test split $0.551\to0.809$, SelCom $0.220\to0.612$; mean across drone surfaces $0.733\to0.834$) and dissolves two artefacts of the reject class at once: the grayscale recall \emph{hole} (below) and the per-class threshold $\tau$ that patched it. The cost is confuser fire the reject net no longer compounds away: RGB-confuser frame fire $0.0015\to0.011$ (still near $1\%$, and damped by temporal voting, Section~\ref{sec:temporal_results}), thermal $0.217\to0.237$, real-video-confuser $0.073\to0.225$. We ship the no-reject router and keep \texttt{robust8} as the ablation that quantifies the reject net; on confuser-rich \emph{paired} streams (the Svanstr\"om flagship), where both modalities are live and a \texttt{reject} is rarely wrong, \texttt{robust8} remains the stronger choice (at the shipped \texttt{filt$\to$clf} order, $0.963$ vs $0.944$) and is the recommended configuration there. (The composed-F1 recovery figures above compare \texttt{robust8} and \texttt{robust8-nr} at a fixed \texttt{clf$\to$filt} composition, isolating the reject variable; the shipped \texttt{filt$\to$clf} order, chosen in Section~\ref{sec:pipeline_paired}, lifts the no-reject numbers further still, e.g.\ Svanstr\"om $0.930\to0.944$ and DUT $0.792\to0.837$.)

\paragraph{The grayscale hole, and why the reject class caused it.}
\texttt{robust6} inherits one honest weakness from the four-class formulation: in the grayscale-IR fallback the thermal features go uninformative, so drone-only-in-RGB frames are routed to \texttt{reject} and grayscale recall collapses. \texttt{robust8} patched this \emph{within} the reject formulation by adding \texttt{rgb\_mean\_conf} and a per-class threshold $\tau$ on $P(\text{trust\_rgb})$: at $\tau=0.20$ it lifts Svanstr\"om-grayscale recall $0.577 \to 0.681$ (Figure~\ref{fig:robust8_operating}), at the structural cost of $7$--$10\times$ real-video-confuser fire on the grayscale slice ($0.006 \to 0.046$), because \texttt{rgb\_mean\_conf} reads as a drone on Svanstr\"om but as a confident bird on real-motion video. Dropping the reject class removes the hole at its source: a router that never rejects cannot route a grayscale drone to \texttt{reject}, so \texttt{robust8-nr} needs neither the $\tau$ dial nor the precision/recall trade the figure illustrates. The sweep is retained as the design-history record of why the reject formulation was awkward to operate.
% [source: ledger=robust8-grayscale-router; eval=svan_gray_robust8_t20_clffilter,video_confuser_robust8_t20_clffilter; run=eval/compare_routing_pipeline.py; cache=eval/results/_routing_pipeline_cmp/comparison.md]

\begin{figure}[htbp]
    \centering
    \includegraphics[width=0.7\textwidth]{fig_robust8_operating_point}
    \caption{Design-history: the grayscale \texttt{trust\_rgb} hole was a decision-rule problem inside the four-class router, not a feature-space limit. Sweeping the per-class \texttt{trust\_rgb} threshold $\tau$ trades precision for recall along a well-separated curve; the reject-class \texttt{robust8} operated at $\tau{=}0.20$. The shipped \texttt{robust8-nr} drops the reject class and removes this knob entirely (the hole cannot form without a \texttt{reject} label). Source: \texttt{eval/compare\_routing\_pipeline.py}.}
    \label{fig:robust8_operating}
\end{figure}

\paragraph{A dual-classifier extension (validated, not shipped).}
Splitting the router into a paired model and a separate grayscale model strictly dominates single \texttt{sa32} on missing-modality cases (RGB-fallback bird-confuser drone $F1$ $0.247\to1.000$; \texttt{rgb\_dataset\_test} $0.690\to0.960$) with paired-core performance unchanged; grayscale airplane/helicopter cases remain weak ($F1\approx0.50$). It is recorded as a validated extension rather than the shipped default; \texttt{robust8}'s two extra features close most of the same gap without a second model.
% [source: ledger=dual-classifier-v3; eval=dualclf_v3_vs_sa32; config=dual_clf_v3_surfaces]

\paragraph{Design evolution: the OOD confuser zoo (predecessor configuration).}
The first quantitative router comparison ran the predecessor cascade (baseline RGB, alert-gated patch v2, \texttt{imgsz=1280}) over the confuser zoo. The open-world \texttt{fusion\_no\_fn\_v1.1} suppressed the bare $52.1\%$ fire to $1.6\%$ at the classifier stage ($0.8\%$ after the filter); \texttt{sa32} reached only $20.5\%$ ($10.3\%$), a $13\times$ gap on this surface.
% [source: ledger=cascade-confuser-collapse; eval=clfzoo_fnfn,clfzoo_sa32; run=eval/cumulative_halluc.py; cache=eval/results/_cumulative_halluc/confuser_fusion_no_fn_model_v1.1/summary.json; config=confuser_zoo_1280 (predecessor configuration)]

\begin{table}[htbp]
    \centering
    \caption{Design-evolution record: classifier comparison under the predecessor configuration (baseline RGB, v2 patch filter, \texttt{imgsz=1280}; Svanstr\"om stride 9). The split verdict between Svanstr\"om and the OOD zoo is the point.}
    \label{tab:classifiers}
    \begin{tabular}{lccccccc}
        \toprule
        Classifier & thr & Drone S2 $R$ & Drone S2 F1 & Drone S3 $R$ & Drone S3 F1 & Zoo S2 & Zoo S3 \\
        \midrule
        \texttt{fusion\_no\_fn\_v1.1}        & 0.8/0.9 & 0.909 & 0.912 & 0.856/0.869 & 0.889/0.895 & \textbf{0.016} & \textbf{0.008} \\
        \textbf{\texttt{scene\_aware\_v3more\_32feat}} & 0.8 & 0.922 & 0.919 & 0.868 & 0.896 & 0.205 & 0.103 \\
        \texttt{control\_v3more\_40feat} & 0.9 & \textbf{0.934} & \textbf{0.925} & \textbf{0.893} & \textbf{0.909} & 0.212 & 0.094 \\
        \bottomrule
    \end{tabular}
\end{table}

\begin{figure}[htbp]
    \centering
    \includegraphics[width=0.7\textwidth]{fig8_classifier_reversal}
    \caption{No trust classifier wins both surfaces. On the OOD confuser zoo \texttt{fusion\_no\_fn\_v1.1} is far the safest (S2 fire $0.016$) but on real video it collapses drone $F1$ to $0.219$; \texttt{sa32} inverts the ranking, dominating real-video drone $F1$ ($0.826$) at a higher zoo fire ($0.205$). This surface reversal motivated the \texttt{sa32} working pick, later superseded by the statistically-derived \texttt{robust6}/\texttt{robust8} (Section~\ref{sec:feature_selection}) and the shipped no-reject \texttt{robust8-nr} (Section~\ref{sec:classifier_results}).}
    % [source: ledger=three-classifier-realvideo; eval=pipe3clf_sa32,pipe3clf_fnfn,clfzoo_fnfn; config=pipe_video_drone_iop]
    \label{fig:classifier_reversal}
\end{figure}

\paragraph{The surface reversal.}
On real video the ranking inverted: \texttt{sa32} dominated segment drone $F1$ ($0.826$ baseline RGB) while \texttt{fusion\_no\_fn\_v1.1} accepted only about one in eight of the correct RGB drone TPs (segment $R=0.128$; seven of eight drone events missed) and \texttt{control\_v3more\_40feat} trailed \texttt{sa32} by 18~pp with no compensating win anywhere. The Svanstr\"om/zoo ordering does not select the production classifier; the operational surface does. \texttt{fusion\_no\_fn\_v1.1} is retained as the open-world fallback for deployments where missed drones cost less than false alarms. This reversal is what discredited hand-ranking classifiers on single benchmarks and motivated the statistical selection that produced \texttt{robust6}/\texttt{robust8} above.
% [source: ledger=three-classifier-realvideo; eval=pipe3clf_sa32,pipe3clf_fnfn; config=pipe_video_drone_iop]

\subsection{Confuser Filter}
\label{sec:verifier_results}
\label{sec:patch_threshold}
\label{sec:patch_audit}

\paragraph{The patch generation: useful, gated, and bounded.}
The patch filter's operating point was set by a threshold sweep on Svanstr\"om (stride 9, 3{,}190 frames): at \texttt{patch\_thr}$=$0.9 it catches $\sim$63\% of the classifier stage's residual confuser FPs at $-4.0$~pp drone F1, the production choice of its era; at $0.5$ the recall cost doubles. Its veto distribution is bimodal (confusers $\gtrsim 0.9$, drone TPs near $0$), which is what makes the threshold a usable knob.
% [source: ledger=patch-catch-below-bar; eval=svan_s3_sa32_thr08; config=svan_iop_1280_s9]

\begin{table}[htbp]
    \centering
    \caption{Patch v2 catch and veto rates by bucket at \texttt{patch\_thr=0.5} (Svanstr\"om @1280, stride 9, 3{,}190 frames, 3{,}130 detections). All three confuser buckets sit below the $0.90$ decisiveness bar.}
    \label{tab:patch_audit}
    \begin{tabular}{lrcr}
        \toprule
        Bucket & N & Catch / veto rate & Median patch prob \\
        \midrule
        BIRD       & 807 & 0.638 catch & 0.904 \\
        AIRPLANE   & 532 & 0.517 catch & 0.540 \\
        HELICOPTER & 464 & 0.709 catch & 0.987 \\
        DRONE\_TP (unwanted veto) & 1248 & 0.054 veto & 0.000 \\
        DRONE\_FP (bonus catches) &   79 & 0.266 catch & 0.000 \\
        \bottomrule
    \end{tabular}
\end{table}
% [source: ledger=patch-catch-below-bar; eval=patch_catch_v2_svan; config=patch_catch_svan_1280_s9]

\begin{figure}[htbp]
    \centering
    \includegraphics[width=0.8\textwidth]{fig8_patch_catchbar}
    \caption{Patch-filter (v2) per-bucket confuser catch rate at \texttt{patch\_thr}=0.5, against the $0.90$ ``decisiveness'' bar. Helicopters ($71\%$) and birds ($64\%$) are mid; airplanes ($52\%$, median patch probability $0.540$) are the actionable gap. All three buckets sit below the bar, the motivation for the feature-reuse successor.}
    % [source: ledger=patch-catch-below-bar; eval=patch_catch_v2_svan; config=patch_catch_svan_1280_s9]
    \label{fig:patch_catchbar}
\end{figure}

The audit's verdict (every confuser bucket below the decisiveness bar, airplanes worst, and a hard distribution bound: on SelCom-scale crops it is bit-for-bit the bare detector) is what motivated replacing pixel re-processing with feature reuse.
% [source: ledger=patch-catch-below-bar,patch-v2-neutral-selcom; eval=patch_catch_v2_svan,v5_selcom_patch]

\paragraph{The feature-reuse generation.}
\begin{table}[htbp]
    \centering
    \caption{Feature-reuse filter (\texttt{mlp\_v5}) vs the MobileNetV3 patch filter (v2), both on the FT4 RGB detector. $F1$ is drone-class (IoP@0.5 on Svanstr\"om/SelCom, IoU@0.5 on Anti-UAV/\texttt{rgb\_dataset}); \emph{halluc} is the confuser hallucination rate. \texttt{mlp\_v5} wins every confuser-rich surface and regresses on the photo-style \texttt{rgb\_dataset} split: the carve-out.}
    \label{tab:distill_verifier}
    \begin{tabular}{lcccccc}
        \toprule
        & \multicolumn{3}{c}{Drone $F1$} & \multicolumn{3}{c}{Halluc. rate} \\
        \cmidrule(lr){2-4} \cmidrule(lr){5-7}
        Surface & bare FT4 & + patch v2 & + \texttt{mlp\_v5} & bare & patch v2 & \texttt{mlp\_v5} \\
        \midrule
        Svanstr\"om (IoP@0.5)      & 0.596 & 0.768 & \textbf{0.869} & 0.470 & 0.163 & \textbf{0.037} \\
        Anti-UAV (IoU@0.5)         & 0.986 & 0.986 & 0.985          & 0.011 & 0.011 & \textbf{0.010} \\
        SelCom CCTV (IoP@0.5)      & 0.591 & 0.591 & \textbf{0.607} & 0.071 & 0.071 & \textbf{0.019} \\
        \texttt{rgb\_dataset} (IoU)& \textbf{0.929} & 0.904 & 0.792 & 0.028 & 0.026 & \textbf{0.010} \\
        confuser-only              & ---   & ---   & ---            & 0.317 & 0.107 & \textbf{0.008} \\
        \bottomrule
    \end{tabular}
\end{table}
% [source: ledger=v5-beats-patch,v5-rgbds-ceiling; eval=v5_svan_mlp,v5_rgbds_mlp,v5_confuser_mlp; run=eval/eval_v4_vs_patch.py; config=svan_iop_1280_s9,rgb_dataset_iou_640,confuser_test_640]

\begin{figure}[htbp]
    \centering
    \includegraphics[width=0.95\textwidth]{fig8_distill_verifier}
    \caption{Feature-reuse filter (\texttt{mlp\_v5}, green) vs the MobileNetV3 patch filter (v2, blue) and the bare FT4 detector (grey), per evaluation surface. (a) Drone $F1$: \texttt{mlp\_v5} wins Svanstr\"om and SelCom, ties saturated Anti-UAV, and regresses on the photo-style \texttt{rgb\_dataset} split. (b) Confuser hallucination rate: lowest on every surface. Generated from \texttt{knowledge/evals.csv}.}
    % [source: ledger=v5-beats-patch; eval=v5_svan_mlp; run=eval/eval_v4_vs_patch.py; config=svan_iop_1280_s9]
    \label{fig:distill_verifier_bar}
\end{figure}

\paragraph{The carve-out, diagnosed.}
\label{sec:mlp_recall_drop}
On the photo-style \texttt{rgb\_dataset} split the filter regresses to $F1=0.792$ ($-11$~pp), driven by a drone recall collapse ($0.896 \to 0.664$). The Model MRI localises why: the falsely-vetoed drones are not low-confidence (confidence-scalar difference between vetoed and kept: $0.000$). They are \emph{smaller}, with a distinct deep-feature signature, and they sit \emph{farther from the confuser distribution than the kept drones do} (centroid distance $16.5$ vs $11.1$) while being far from the filter's training-drone cluster ($15.4$). The filter rejects them because they do not resemble its training drones, not because they resemble birds: an OOD \emph{coverage} gap, not a structural overlap.
% [source: ledger=mlp-v5-recall-drop-is-ood-coverage; run=eval/diagnose_mlp_recall_drop.py; cache=eval/results/_offline_pipeline/cache]

A no-retrain fail-open gate (abstain when a detection is far from the confuser reference; $k$-NN distance) recovers $91$--$100\%$ of the falsely-vetoed drones at a $5$--$10\%$ tolerated confuser leak, and is \emph{not adopted}: calibrated on photo-style confusers it backfires on cluttered Svanstr\"om (precision $0.887 \to 0.631$, because Svanstr\"om's clutter is itself OOD to the calibration set and gets wrongly released). Re-weighting is a no-op (the OOD drones are absent from training) and temporal voting cannot recover a systematic veto. The carve-out is therefore retained rather than patched: on video, frame-level alert recall with the full-veto filter is $\approx$1.0 (a drone seen across consecutive frames survives individual vetoes), so the regression is confined to single-frame photo-style benchmarks, and the patch filter remains the documented fallback on that one surface. The frontier-moving remedy, closing the coverage gap with diverse small-drone data, is the same lever that made the IR filter recall-safe.
% [source: ledger=filter-recall-precision-decision; run=eval/test_failopen_verifier.py,eval/eval_failopen_prepost.py,eval/failopen_expanded_ref.py; cache=eval/results/_offline_pipeline/cache]

\begin{figure}[htbp]
    \centering
    \includegraphics[width=0.7\textwidth]{fig8_failopen_expanded}
    \caption{Why fail-open is not adopted: the recall/precision frontier on Svanstr\"om for the fail-open gate with the original \texttt{rgb\_confusers} reference (red) versus a reference expanded with held-out Svanstr\"om clutter (green), against the full-veto operating point (star) and the bare detector (square). The full-veto point dominates at high recall; the trade is partly calibration and partly fundamental. Source: \texttt{eval/failopen\_expanded\_ref.py}.}
    \label{fig:failopen_expanded}
\end{figure}

\paragraph{The cross-modal IR filter.}
\label{sec:grayscale_verifier}
The aligned filter (\texttt{mlp\_v5\_ir\_aligned}, Section~\ref{sec:ir_xmodal_verifier}) is evaluated on a held-out gate: CBAM, a novel thermal aerial-confuser set excluded from its training.

\begin{table}[htbp]
    \centering
    \caption{Thermal-deploy filter (\texttt{mlp\_v5\_ir\_aligned}, thermal scaler, \texttt{conf}$=$0.40, \texttt{thr}$=$0.05) vs the bare IR detector; FP in parentheses. CBAM is held out of training. The filter cuts held-out novel-confuser FPs by $69\%$ while preserving thermal drone recall in-distribution ($\Delta R$ at most $-0.007$); the MobileNetV3 patch filter on the same held-out CBAM cuts only $7$ FP.}
    \label{tab:ir_aligned}
    \begin{tabular}{lrccc}
        \toprule
        Surface & $n$ & bare IR (\texttt{v3b}) $P/R/F1$ (FP) & + aligned MLP $P/R/F1$ (FP) & $\Delta R / \Delta F1$ \\
        \midrule
        \textbf{CBAM (held out)}  & 180  & $0.547/0.967/0.699$ (48) & $\mathbf{0.786/0.917/0.846}$ (\textbf{15}) & $-0.050 / \mathbf{+0.147}$ \\
        \texttt{ir\_dset\_final}  & 4806 & $0.965/0.965/0.965$ (109) & $0.965/0.958/0.962$ (108) & $-0.007 / -0.003$ \\
        \texttt{ir\_video} test   & 831  & $0.909/0.977/0.942$ (80) & $0.909/0.977/0.942$ (80) & $0.000 / 0.000$ \\
        Anti-UAV test             & 4269 & $0.983/0.942/0.962$ (68) & $0.983/0.942/0.962$ (68) & $0.000 / 0.000$ \\
        \midrule
        \emph{patch filter on CBAM} & 180 & --- & $0.564/0.883/0.688$ (41) & (cut only 7 FP) \\
        \bottomrule
    \end{tabular}
\end{table}
% [source: ledger=ir-grayscale-harvest-solves-thermal-filter; eval=ir_aligned_cbam_heldout,ir_aligned_irdset_final; run=mri_train_aligned,eval_run_aligned_full; config=cbam_ir_640,ir_final_640]

Two properties define this filter: a drone-diversity re-mine ($\sim$30k corrective-set thermal drones) keeps it recall-safe on held-out thermal, and the grayscale-harvested confusers, aligned into thermal feature space, provide its confuser-rejection signal (Section~\ref{sec:ir_xmodal_verifier}). The same network's grayscale-scaler behaviour (a $96.8\%$ FP cut on grayscale confusers) is reported with the grayscale finding (Section~\ref{sec:grayscale}), where its drone-recall bound on that surface is also stated plainly. Across both modalities the feature-reuse generation beats the patch generation on every confuser-rich surface measured; the airplane category remains its weakest, inherited from training distribution rather than architecture.
% [source: ledger=ir-recall-fixed-by-drone-diversity,mlp-beats-patch-both-modalities]

\paragraph{Operating points: the recall/false-positive dial.}
Each filter's threshold is a tunable operating point, not a fixed property of the network. Figure~\ref{fig:filter_operating} sweeps the per-detection $P(\text{drone})$ threshold of all three filters, plotting drone-detection recall (the fraction of true-drone detections retained, pooled over that filter's drone surfaces) against the confuser fire-rate on the matching confuser surface. Three shapes emerge. The RGB \texttt{mlp\_v5} sits on a flat high-recall shoulder: at the shipped $0.25$ it keeps $95.6\%$ of true-drone detections at a $1.1\%$ confuser fire-rate, and stays recall-transparent ($\ge 98\%$) well below that, so its recall is essentially free to recover. The IR-thermal \texttt{aligned} filter is flat in fire at \emph{every} threshold (fire-rate $\approx 0.21$--$0.24$ throughout): the thermal airplanes it misses are a coverage limit (Section~\ref{sec:limits}), not a threshold one, so no setting suppresses them. The grayscale scaler has the steepest recall slope: at the shipped $0.25$ it retains only $46.7\%$ of grayscale-fallback drone detections, the documented cost of the cross-domain mode. The reading for a deployer is that the shipped points sit deliberately on the high-precision side. Where the limit is the threshold (RGB), recall is recoverable almost for free by lowering it; where the limit is coverage (the IR-thermal airplanes, the grayscale cross-domain drones), lowering helps only partially and at a steeper false-positive cost: the grayscale scaler at $0.05$ instead of $0.25$, for instance, lifts detection recall $0.47 \to 0.62$ while its confuser fire-rate rises $0.008 \to 0.036$. The shipped configuration chooses precision where confusers dominate; a recall-first deployment moves left along these curves.
% [source: run=eval/filter_operating_sweep.py; cache=eval/results/filter_operating_sweep.json; surfaces per tab:ablation_confusers]

\begin{figure}[htbp]
    \centering
    \includegraphics[width=\textwidth]{fig_filter_operating}
    \caption{Filter operating points (detection level). For each confuser filter, drone detection recall (blue, left axis; fraction of GT-matched true-drone detections retained, pooled over the drone surfaces) and confuser fire-rate (red dashed, right axis; matching confuser surface) as the per-detection $P(\text{drone})$ threshold varies; the shipped point ($t=0.25$ RGB, $0.05$ IR-thermal, $0.25$ grayscale) is marked. RGB sits on a flat high-recall shoulder; the IR-thermal fire-rate is threshold-insensitive (the airplane coverage gap); the grayscale scaler trades recall steeply (the cross-domain fallback). Source: \texttt{eval/filter\_operating\_sweep.py}.}
    \label{fig:filter_operating}
\end{figure}

% ======================================================================
\section{Model MRI Findings}
\label{sec:mri_findings}

The instrument is described in Section~\ref{sec:model_mri}; this section reports what it found. The through-line: the discrimination the pipeline needs already exists inside the detectors' feature spaces, linearly accessible, in both modalities. That is why the production filters are 517-D MLPs rather than CNNs, and why the trust classifier needs eight features rather than thirty-two.

\begin{figure}[htbp]
    \centering
    \begin{subfigure}{0.6\textwidth}\includegraphics[width=\textwidth]{fig8_mri_lda}\caption{}\end{subfigure}
    \begin{subfigure}{0.39\textwidth}\includegraphics[width=\textwidth]{fig8_mri_anova}\caption{}\end{subfigure}
    \caption{Model MRI of the FT4 detector's fused \texttt{p3}+\texttt{p5} ROI features on the shipped $32{,}931$-detection corpus. (a) The linear-discriminant projection separates drone (blue) from confuser (red) into two cleanly-resolved peaks at $\approx 95\%$ single-threshold accuracy. (b) Per-block ANOVA $F$-statistic: \texttt{p5} is the dominant discriminative layer (max $F=42{,}346$, mean $2{,}006$), ahead of \texttt{p3} and the detection metadata; the confidence scalar ranks only 6th.}
    % [source: ledger=mri-v5-report-regen,v5-lda-separability; eval=distill_cv_mlp; run=mri/cli.py; cache=mri/results/v5_report_regen/stats.json; config=distill_cv_5fold]
    \label{fig:mri_stats}
\end{figure}

\paragraph{The signal is linear, supervised, and deep.}
On the RGB detector, a single Fisher axis separates drone from confuser at $\approx$95\% threshold accuracy while unsupervised PCA finds near-zero cluster structure (silhouette $0.067$): the signal is real but invisible to distance-based methods, which is why a small trained MLP is the right reader. The strongest single discriminator is a \texttt{p5} semantic channel ($F=42{,}346$), $4\times$ stronger than the detection-confidence scalar: the filter reads \emph{what the detector saw}, not \emph{how sure it was}. The IR detector is the thermal twin: LDA $0.981$ over $14{,}697$ drone / $1{,}386$ confuser detections, strongest neuron $F=5{,}370$, top discriminators again \texttt{p5} channels with \texttt{meta:conf} far down the ranking.
% [source: ledger=v5-lda-separability,ir-features-separable; run=mri/cli.py; cache=mri/results/v5_report_regen/stats.json,mri/results/ir_v3b_report]

\begin{figure}[htbp]
    \centering
    \begin{subfigure}{\textwidth}\includegraphics[width=\textwidth]{fig8_mri_act_drone}\caption{Drone (Svanstr\"om, conf$=0.72$)}\end{subfigure}
    \\[4pt]
    \begin{subfigure}{\textwidth}\includegraphics[width=\textwidth]{fig8_mri_act_confuser}\caption{Confuser (\texttt{rgb\_confusers} test, conf$=0.29$)}\end{subfigure}
    \caption{Model MRI spatial ``brain scan'' of the top discriminative neurons. Each row: the detection crop, the \texttt{p3} activation (stride~8, fine spatial detail), and the \texttt{p5} activation (stride~32, semantic depth). On the drone (a) \texttt{p3} binds to the airframe and \texttt{p5} forms a centred semantic blob; on the confuser (b) the same neurons fire diffusely across background structure. This activation-signature difference is what the feature-reuse filter reads.}
    \label{fig:mri_activation}
\end{figure}

\paragraph{Cross-modal alignment is an affine offset.}
The grayscale$\leftrightarrow$thermal comparison (Section~\ref{sec:ir_xmodal_verifier}) is the MRI's most consequential measurement: per-modality z-scoring lifts grayscale$\to$thermal confuser-transfer from chance to near the within-thermal ceiling, far above full-covariance CORAL (Table~\ref{tab:gray_thermal_auroc}): the modality gap in the detector's feature space is a diagonal mean/scale shift, not a representational change. This is what lets the abundant grayscale-confuser corpus train directly into thermal feature space.
% [source: ledger=gray-thermal-alignable; run=mri/modality_align.py; cache=mri/results/ir_v3b_report; config=ir_v3b]

\begin{figure}[htbp]
    \centering
    \includegraphics[width=\textwidth]{fig9_ir_gray_align}
    \caption{Grayscale$\leftrightarrow$thermal feature alignment in the IR detector. (A) After a per-modality z-score, thermal and grayscale drones land in the same region of the drone-vs-confuser discriminant axis, and likewise for confusers: the modalities are interchangeable up to an affine offset. (B) Grayscale-trained confuser-filter transfer to thermal rises from chance ($0.500$) to near the within-thermal ceiling ($0.919$ vs $0.974$) once the offset is removed; CORAL ($0.707$) under-performs the simpler z-score.}
    % [source: ledger=gray-thermal-alignable; run=mri/modality_align.py; cache=mri/results/ir_v3b_report; config=ir_v3b]
    \label{fig:ir_gray_align}
\end{figure}

\paragraph{The instrument audits its own paper trail.}
Re-running the MRI on the shipped V5 corpus corrected two stale figures and one headline claim in an earlier hand-written report (the dominant discriminator is \texttt{p5}, not the confidence score). The same diagnostic pattern (localise the failure in feature space, prescribe the recall-safe fix) recovered both the RGB carve-out diagnosis (Section~\ref{sec:verifier_results}) and the IR filter (Section~\ref{sec:ir_xmodal_verifier}).
% [source: ledger=mri-v5-report-regen; run=mri/cli.py; cache=mri/results/v5_report_regen/stats.json]

% ======================================================================
\section{The Grayscale Finding}
\label{sec:grayscale}

The IR detector was trained exclusively on thermal infrared data. The grayscale-RGB mode (Section~\ref{sec:ir_grayscale_mode}) exists as a hardware fallback: convert the RGB stream to single-channel grayscale, feed it to the thermal model. It was built as a conservative degradation mode; it turned out to be competitive, and that result, plus the harvesting path it opened, is the finding.

\subsection{The Fair Three-Way Comparison}

On Svanstr\"om, at each model's canonical configuration and against the same ground truth:

\begin{table}[htbp]
    \centering
    \caption{The grayscale three-way on Svanstr\"om ($n=4{,}000$ of 28{,}710, IoP@0.5, bare detectors). The thermal-trained detector on grayscale-converted RGB performs within $2.7$~pp F1 of the dedicated RGB detector (with higher precision and lower recall), while the raw-RGB control isolates the grayscale conversion as the load-bearing step.}
    \label{tab:gray_threeway}
    \begin{tabular}{lccc}
        \toprule
        configuration & $P$ & $R$ & F1 [95\% CI] \\
        \midrule
        RGB detector (\texttt{ft4}) on RGB, @1280 & 0.454 & 0.916 & 0.607 [0.590--0.623] \\
        IR detector (\texttt{v3b}) on raw RGB, @640 & 0.362 & 0.126 & 0.187 [0.165--0.210] \\
        IR detector (\texttt{v3b}) on grayscale RGB, @640 & 0.543 & 0.621 & 0.580 [0.561--0.599] \\
        \bottomrule
    \end{tabular}
\end{table}
% [source: thesis_eval/results/tier1_results.json (svanstrom A ft4/rgb; svanstrom_rawrgb A; svanstrom_gray A v3b/ir)]

The control row settles the mechanism question at benchmark grade: fed the raw three-channel RGB, the thermal model collapses to $F1=0.187$; fed the \emph{same frames} through a one-line grayscale conversion, it reaches $0.580$, a tripling that isolates the single-channel intensity input as the load-bearing step. The zero-shot thermal model is then $2.7$~pp F1 behind the dedicated RGB detector on visible-light frames, with the opposite error profile (more precise, less complete). On the real-video diagnostic, where the finding was first observed, the same mode more than doubles the raw-RGB control's aggregate F1 ($0.295 \to 0.636$), nearly halves bird-confuser FPPI ($0.261 \to 0.142$), and on the hardest bird-cluttered clip ties the best RGB detector ($F1$ $0.837$ vs $0.840$, a tie at that clip's sample size):
% [source: ledger=ir-grayscale-fallback; eval=vid_drone_ir_gray,vid_conf_ir_gray; run=eval/eval_video_tests.py; config=video_drone_iop,video_confuser]

\begin{table}[htbp]
    \centering
    \caption{Real-video six-mode diagnostic (the finding's discovery surface; 9 drone clips / 1{,}234 GT instances, 10 confuser clips / 1{,}250 frames, IoP@0.5). The IR-on-grayscale rows are the discovery; the raw-RGB row is the control.}
    \label{tab:realvideo_master}
    \begin{tabular}{lccccccc}
        \toprule
        & \multicolumn{3}{c}{Drone (9 videos)} & \multicolumn{4}{c}{Confuser FPPI (10 videos)} \\
        \cmidrule(lr){2-4} \cmidrule(lr){5-8}
        Mode & $P$ & $R$ & $F1$ & All & Bird & Airplane & Helicopter \\
        \midrule
        RGB baseline           & 0.771 & 0.749 & \textbf{0.760} & 0.512 & 0.960 & 0.451 & 0.278 \\
        RGB \texttt{retrained\_v2} & \textbf{0.909} & 0.453 & 0.605 & 0.196 & 0.139 & 0.385 & 0.133 \\
        RGB \texttt{selcom\_1280}  & 0.649 & \textbf{0.812} & 0.721 & 0.709 & 1.528 & 0.493 & 0.333 \\
        RGB \texttt{selcom\_640}   & 0.807 & 0.666 & 0.730 & 0.260 & 0.313 & 0.359 & 0.179 \\
        \textbf{IR on grayscale-RGB} & 0.743 & 0.557 & 0.636 & \textbf{0.158} & \textbf{0.142} & \textbf{0.283} & \textbf{0.103} \\
        IR on raw RGB          & 0.647 & 0.191 & 0.295 & 0.150 & 0.261 & 0.158 & 0.079 \\
        \bottomrule
    \end{tabular}
\end{table}
% [source: ledger=ir-grayscale-fallback,no-pareto-dominance-video; eval=vid_drone_*,vid_conf_*; run=eval/eval_video_tests.py; config=video_drone_iop]

\begin{figure}[htbp]
    \centering
    \includegraphics[width=0.95\textwidth]{fig_grayscale_panel}
    \caption{Cross-modal transfer, illustrated on one Anti-UAV night frame (each panel cropped around the drone). \emph{Left}: the production RGB detector (\texttt{ft4}) on the RGB frame, for reference. \emph{Middle}: the IR detector (\texttt{v3b}), trained zero-shot on thermal only, applied to the grayscale-converted frame; it lands the drone (IoU $0.81$ with the green GT box). \emph{Right}: the same IR detector on the raw 3-channel RGB collapses (IoU $0.00$, no box). The transfer is carried by the single-channel intensity input; the colour information that distinguishes thermal from visible is what breaks it.}
    \label{fig:grayscale_qualitative}
\end{figure}

\subsection{Mechanism and Bounds}

The likely mechanism is shared low-level statistics: thermal IR and grayscale RGB are both single-channel intensity images that render small drones as silhouettes against sky, and a detector trained on thermal silhouette gradients retains them when colour, the one channel property that genuinely differs, is removed (Figure~\ref{fig:grayscale_qualitative}). The transfer is partial and patterned: on clean takeoff clips the dedicated RGB detectors win by $28$--$41$~pp F1 per clip (texture-rich close-ups favour visible-light training), while on small-silhouette bird-cluttered scenes the thermal model's discrimination dominates. The IR model has also never seen a visible-light bird, yet its grayscale bird-FPPI ($0.142$) lands within sampling error of a detector explicitly hard-negative-mined for them ($0.139$): thermal-domain confuser training transfers along with the detection.
% [source: ledger=ir-grayscale-fallback; eval=vid_conf_ir_gray; config=video_confuser]

To our knowledge, no published study quantifies this specific transfer: a thermal-only-trained drone detector evaluated zero-shot on grayscale-converted visible-light video against drone ground truth. The adjacent literatures (modality hallucination, RGB-to-thermal translation, domain-adversarial alignment) all provide cross-modal information at training time (Section~\ref{sec:lit_xmodal_cascade}); here the transfer emerges from shared low-level structure alone.

\subsection{The Harvesting Payoff and the Honest Bound}

The finding's practical value is double. As a \emph{fallback}, grayscale mode gives a thermal-equipped system a usable second channel when the thermal camera is absent. As a \emph{harvesting surface}, it is where the thermal filter's confuser training comes from: the IR detector hallucinates $\sim$20$\times$ more on grayscale confusers than on thermal ones ($37.2\%$ vs $1.8\%$ per image), so grayscale yields abundant confuser detections, and the MRI's alignment result (Section~\ref{sec:mri_findings}) carries them into thermal feature space, yielding \texttt{mlp\_v5\_ir\_aligned} (Section~\ref{sec:verifier_results}).
% [source: ledger=ir-grayscale-is-hallucination-mode; run=mri/modality_probe.py; cache=mri/results/ir_v3b_report]

On the grayscale channel itself, the aligned filter (grayscale scaler) cuts confuser FPs by $96.8\%$ ($656 \to 21$ on the grayscale-converted confuser corpus), and over-vetoes drones on grayscale Svanstr\"om: a threshold sweep over the cached probabilities finds no operating point that preserves useful drone recall on that surface (recall $\leq 0.27$ even at threshold $0.02$), because the filter's grayscale drone training came from photo-style imagery, not Svanstr\"om-scale silhouettes. The deployment consequence is stated plainly: on the grayscale path the filter is a \emph{confuser-suppression} tool, gated rather than unconditional; grayscale \emph{detection} runs unfiltered ($F1=0.580$ bare). Grayscale is a fallback and a harvesting surface, not a primary detection surface.
% [source: thesis_eval/results/tier1_results.json (gray_confuser C, svanstrom_gray S4 + GRAY_SWEEP)]

% ======================================================================
\section{Threats to Validity}
\label{sec:threats}

\paragraph{Internal validity.}
Three threats. First, the RGB detector's training stances and the patch filter's versions were iterated against Svanstr\"om evaluation results; some design decisions are implicitly biased toward Svanstr\"om performance. The real-video clips and the SelCom footage, neither consulted during those iterations, are the partial mitigation. Second, \texttt{imgsz} acts as a hidden hyperparameter; the protocol (Section~\ref{sec:canonical_configs}) pins it per surface, and the one cross-resolution comparison the thesis needs (Svanstr\"om RGB@1280 vs IR@640) is flagged where it appears. Third, the trust classifier was trained against a specific RGB confidence calibration; swapping the RGB detector without re-mining classifier data risks silent degradation (the documented \texttt{retrained\_v2\_32feat} failure, Section~\ref{sec:classifier_variants}).

\paragraph{Construct validity.}
IoP at 0.5 is more lenient than IoU at 0.5 for small-object detection; numbers here are not directly comparable to IoU-based literature without conversion (Section~\ref{sec:metrics}). Trust-aware scoring is the second construct choice, justified in Section~\ref{sec:scoring_audit}; both rules are declared per table.

\paragraph{External validity.}
The IR detector and the trust classifier train on frames and sequences that overlap both paired evaluation surfaces; the overlap is quantified per component, and bounded by the held-out clean split, in Section~\ref{sec:svanstrom_audit}: Anti-UAV conclusions are unchanged on $57{,}542$ held-out frames ($0.973 \to 0.977$ bare, $0.984 \to 0.986$ pipeline), the Svanstr\"om IR detector's solo numbers carry at most $7.3$~pp of in-distribution inflation (part of it sequence difficulty, per the leakage-free RGB control), and the production cascade moves only $-1.4$~pp ($0.949 \to 0.935$). The deployment evidence is one site (SelCom) and one camera class; generalisation to other CCTV installations is untested.
% [source: runs/clean_split/clean_split_results.json]

\paragraph{Sampling and uncertainty.}
Large surfaces are evaluated on $\approx$4{,}000-frame even-spread samples (Section~\ref{sec:canonical_configs}); all tables print $n$ and the source size, and headline cells carry 95\% bootstrap CIs. Differences within overlapping CIs are treated as ties; this includes the grayscale three-way's $2.7$~pp gap, which is reported as ``slightly behind'', not as parity. Per-clip real-video numbers ($20$--$340$ frames per clip) are indicative; aggregates are the reliable quantities.

\paragraph{Conclusion validity.}
Single-run point estimates are never compared at sub-CI granularity; the production-stack choices are justified on cost and OOD behaviour where in-domain differences are within noise (Section~\ref{sec:pipeline_ablation}).

\section{Limitations and Future Work}
\label{sec:limits}

\paragraph{Edge latency.} All evaluation ran on a GTX 1050~Ti; the cascade is interactive there, but Jetson-class edge latency is unmeasured. The per-stage budget (Table~\ref{tab:speed}) makes the port plausible; it remains future work.

\paragraph{GUI integration of the aligned filter.} The production design runs \texttt{mlp\_v5\_ir\_aligned} per-frame on the IR branch; it is validated in the evaluation engine, while its operator-GUI wiring is an open engineering item at the time of writing.

\paragraph{The airplane gap.} Every filter generation suppresses airplanes worst (patch: $52\%$ catch on Svanstr\"om airplanes, below the decisiveness bar; the feature-reuse filters inherit the gap, and the thermal-confuser surface's residual fire is airplane-dominated). Remediation is data, not architecture: airplane crops from new domains as confuser-class training data.

\paragraph{IR OOD recall.} The IR detector's recall on independent, verified thermal drone footage outside its corpus families is unmeasured (the available community thermal sets could not be quality-verified and are not used); this is the evaluation gap the RGB parallel channel hedges against.

\paragraph{Track-level classification.} The temporal smoother maintains per-track state and renders flight paths, but does not yet classify on track dynamics (velocity, persistence). Drones, birds, and aircraft have distinct flight signatures; a track-based classifier is a largely orthogonal evidence channel and the most promising single direction, particularly for the airplane gap that appearance-based stages cannot close.

\paragraph{Open evaluation items.} The thermal-confuser per-category breakdown formerly registered here has since been run (a CPU replay of the existing cache): the \texttt{IR\_confusers} resistance is airplane-driven ($35.2\%$ bare airplane fire vs $12.2\%$ bird and $0\%$ helicopter), confirming the airplane gap as the active front. The video-tuned filter-threshold question is closed by the sweep in Section~\ref{sec:temporal_results}.
% [source: thesis_eval/results/notes_round1_results.json (ir_confusers CAT)]
